\documentclass[fin1, tisk]{fmfdelo}
% \documentclass[fin1, tisk]{fmfdelo}
% Če pobrišete možnost tisk, bodo povezave obarvane,
% na začetku pa ne bo praznih strani po naslovu, …

%%%%%%%%%%%%%%%%%%%%%%%%%%%%%%%%%%%%%%%%%%%%%%%%%%%%%%%%%%%%%%%%%%%%%%%%%%%%%%%
% METAPODATKI
%%%%%%%%%%%%%%%%%%%%%%%%%%%%%%%%%%%%%%%%%%%%%%%%%%%%%%%%%%%%%%%%%%%%%%%%%%%%%%%

% - vaše ime
\avtor{Luka Lavš}

% - naslov dela v slovenščini
\naslov{Mooreovi grafi in posplošeni večkotniki}

% - naslov dela v angleščini
\title{Moore graphs and generalized polygons}

% - ime mentorja/mentorice s polnim nazivom:
%   - doc.~dr.~ime priimek
%   - izr.~prof.~dr.~ime priimek
%   - prof.~dr.~ime priimek
%   za druge variante uporabite ustrezne ukaze
\mentor{doc.~dr.~Janoš Vidali}
% \somentor{...}
% \mentorica{...}
% \somentorica{...}
% \mentorja{...}{...}
% \somentorja{...}{...}
% \mentorici{...}{...}
% \somentorici{...}{...}

% - leto diplome
\letnica{2026} 

% - povzetek v slovenščini
%   V povzetku na kratko opišite vsebinske rezultate dela. Sem ne sodi razlaga
%   organizacije dela, torej v katerem razdelku je kaj, pač pa le opis vsebine.
\povzetek{
    V tej diplomski nalogi bomo predstavili Mooreeve grafe, navedli njihove lastnosti, 
    pokazali povezavo z kletkami, predstavili njihove konstrukcije ter končali s povezavo
    o posplošenih večkotnikih in goemetrijah.
}

% - povzetek v angleščini
\abstract{...}

% - klasifikacijske oznake, ločene z vejicami
%   Oznake, ki opisujejo področje dela, so dostopne na strani https://www.ams.org/msc/
\klasifikacija{..., ...}
%05C50 05C25 20B25 20C15

% - ključne besede, ki nastopajo v delu, ločene s \sep
\kljucnebesede{...\sep ...}

% - angleški prevod ključnih besed
\keywords{...\sep ...} % angleški prevod ključnih besed

% - angleško-slovenski slovar strokovnih izrazov
\slovar{
% \geslo{angleški izraz}{slovenski izraz}
% ...
}

% - ime datoteke z viri (vključno s končnico .bib), če uporabljate BibTeX
\literatura{literatura.bib}

%%%%%%%%%%%%%%%%%%%%%%%%%%%%%%%%%%%%%%%%%%%%%%%%%%%%%%%%%%%%%%%%%%%%%%%%%%%%%%%
% DODATNE DEFINICIJE
%%%%%%%%%%%%%%%%%%%%%%%%%%%%%%%%%%%%%%%%%%%%%%%%%%%%%%%%%%%%%%%%%%%%%%%%%%%%%%%

% naložite dodatne pakete, ki jih potrebujete
% \usepackage{sagetex,xcolor,tikz,tkz-graph,tkz-berge}
\usepackage{mathrsfs}
\usepackage{enumitem}
\usepackage{graphicx}
\usepackage{color}
\usepackage{import}
\usepackage{url}
\usepackage{tikz}
\usetikzlibrary{calc, decorations.pathreplacing}
\usetikzlibrary{positioning, arrows.meta}

% deklarirajte vse matematične operatorje, da jih bo LaTeX pravilno stavil
% \DeclareMathOperator{\...}{...}

% vstavite svoje definicije ...
\newcommand{\red}{\operatorname{red}}
\newcommand{\srg}{\operatorname{srg}}
\newcommand{\diam}{diam}
\newcommand{\N}{\mathbb{N}}
\newcommand{\ZZ}{\mathbb{Z}}
\newcommand{\reach}{\operatorname{Reach}}
\newcommand{\gr}{\mathscr{G}}

\newcommand{\A}{\mathcal{A}}
\newcommand{\II}{\mathbf{I}}
\newcommand{\IIw}[2]{#1 \;\II\; #2}
\newcommand{\IIwd}[2]{#1 \;\II^D\; #2}
\newcommand{\PP}{\mathcal{P}}
\newcommand{\LL}{\mathcal{L}}
\newcommand{\PG}{\mathrm{PG}}
\newcommand{\GF}{\mathrm{GF}}
\newcommand{\GQ}{\mathrm{GQ}}
\newcommand{\rank}{\operatorname{rank}}
\newcommand{\sled}{\operatorname{sl}}
\newcommand{\aut}{\operatorname{Aut}}
\newcommand{\id}{\operatorname{id}}
\newcommand{\kneser}{\operatorname{KG}}
\newcommand{\specter}{\operatorname{Spec}}
\newcommand{\im}{\operatorname{im}}
\newcommand{\sgn}{\operatorname{sgn}}

\tikzset{
    vnode/.style={circle, draw, fill=black, inner sep=0pt, minimum size=3pt},
    old_struct/.style={thin, black!40},
    old_edge/.style={thin, black!40},
    new_edge/.style={very thick, black},
    main_arrow/.style={-{Stealth}, line width=0.3pt, black},
    label_style/.style={font=\tiny, text=black!70}
}

%%%%%%%%%%%%%%%%%%%%%%%%%%%%%%%%%%%%%%%%%%%%%%%%%%%%%%%%%%%%%%%%%%%%%%%%%%%%%%%
% ZAČETEK VSEBINE
%%%%%%%%%%%%%%%%%%%%%%%%%%%%%%%%%%%%%%%%%%%%%%%%%%%%%%%%%%%%%%%%%%%%%%%%%%%%%%%

% https://fran.si/135/epravopis-slovenski-pravopis/4415248/moore?View=1&Query=moore
% 	EDNINA	DVOJINA	MNOŽINA
% imenovalnik	Moore	Moora	Moori
% rodilnik	Moora	Moorov	Moorov
% dajalnik	Mooru	Mooroma	Moorom
% tožilnik	Moora	Moora	Moore
% mestnik	pri Mooru	pri Moorih	pri Moorih
% orodnik	z Moorom	z Mooroma	z Moori

\begin{document}


\section{Uvod v teorijo grofov}



\subsection{Osnovna teorija grafov}

\subsection{Teorija grup}

\begin{definicija}
    Množica generatorjev grupe $G$ je podmnožica $S \subseteq G$, 
    tako da je vsak element $g \in G$ mogoče zapisati kot produkt elementov iz $S$ in njihovih inverzov.
    Pišemo $G = \langle S \rangle$.
\end{definicija}

\begin{definicija}\label{def:normalna_podgrupa}
    Podgrupa $H$ grupe $G$ je normalna podgrupa v $G$, če velja
    $gHg^{-1} = H$ za vsak $g \in G$. Označimo $H \triangleleft G$.
\end{definicija}
\begin{definicija}
    Produkt $HK$ dveh podgrup $H$ in $K$ grupe $G$ je množica $HK = \{hk \mid h \in H, k \in K\}$.
\end{definicija}
\begin{definicija}
    Direktni produkt $G \times H$ dveh grup $(G, \cdot)$ in $(H, \circ)$ je grupa $(G \times H, *)$, kjer so
    $\{(g, h)\mid g\in G, h\in H\}$ njeni elementi, operacija pa definirana kot
    $(g_1, h_1) * (g_2, h_2) = (g_1 \cdot g_2, h_1 \circ h_2)$.
\end{definicija}
\begin{definicija}\label{def:semidirektni_produkt}
    Naj bo $G$ grupa z enoto $e$, $H$ njena podgrupa in $N \triangleleft G$ normalna podgrupa.
    Če je $G$ enak produktu $NH$ in $H \cap N = \{e\}$, potem pravimo, da je $G$ semidirektni produkt $H$ in $N$, 
    kar označujemo z $G = N \rtimes H$. 
\end{definicija}

\begin{opomba}\label{op:lastnost_semidirektnega_produkta}
    V primeru, da je $G = N \rtimes H$, potem je $|G| = |N||H|$, saj je $G = NH$ in $H \cap N = \{e\}$.
\end{opomba}

\begin{trditev}\label{tr:semidirect_order}
    Naj bo $H = \langle\beta\rangle$ grupa generirana z $\beta$. 
    Naj bo $N = \langle \alpha \rangle$ grupa generirana z $\alpha$.
    Grupa $N$ je normalna v $\langle\alpha\beta\rangle$ ($N \trianglelefteq \langle\alpha\beta\rangle$) 
    in $N \cap H = \{e\}$, če velja $\alpha\beta\alpha^{-1} \in \langle \beta \rangle$, oziroma
    $\alpha\beta\alpha^{-1} = \beta^m$ za nek $m \in \mathbb{Z}$.

    Potem je $G = \langle \beta \rangle \rtimes \langle \alpha\rangle$ semidirektni produkt
    in $|G| = |N||H| = |\langle \alpha \rangle||\langle \beta \rangle|$.
\end{trditev}


\begin{definicija}
    Permutacija
\end{definicija}
\begin{definicija}
    Transpozicija je permutacija, ki zamenja le dva elementa, ostale pa pusti na mestu.
    Transpozicijo, ki zamenja elementa $i$ in $j$ ter pusti ostale na mestu, 
    označimo z $\tau: (i\; j)$.
\end{definicija}

\begin{definicija}
    Permutacija $\sigma$ je soda, če se da zapisati kot produkt sodega števila transpozicij.
\end{definicija}

\begin{definicija}
    Simetrična grupa $S_n$ je grupa vseh permutacij množice \\
    $[n] = \{1, 2, \ldots, n\}$, 
    kjer je operacija kompozitum funkcij.
\end{definicija}

\begin{definicija}
    Alternijajoča grupa $A_n$ je podgrupa simetrične grupe $S_n$, ki jo tvorijo vse sode permutacije.
\end{definicija}

\begin{definicija}
    Ciklična grupa $C_n$ je grupa, ki jo tvori element $g$ reda $n$, 
    torej $C_n = \{e, g, g^2, \ldots, g^{n-1}\}$. Alternativno lahko
    ciklično grupo definiramo kot podgrupo simetrične grupe $S_n$, 
    ki jo tvori ciklična permutacija, ki zamenja $n$ elementov v enem ciklu, na primer 
    $(1\; 2\; \ldots \; n)$.
\end{definicija}

\begin{zgled}
    Ciklična grupa $C_2 = \{e, g\}$, kjer je $g^2 = e$. Izomorfna je grupi 
    $\{e, \tau\}$, kjer je $\tau$ transpozicija, na primer $\tau = (1\; 2)$.
\end{zgled}

\begin{zgled}
    Simetrična grupa $S_n$ je semidirektni produkt alternijajoče grupe $A_n$ in cikla $C_2$,
    velja $S_n = A_n \rtimes C_2$,
    ki ga tvorijo soda permutacija. Alternijajoča grupa $A_n$ je 
    normalna podgrupa simetrične grupe $S_n$,
\end{zgled}
\begin{dokaz}
    Alternijajoča grupa $A_n$ je normalna grupa v $S_n$, namreč 
    za vsak $\tau \in A_n$ in $\sigma \in S_n$ je 
    $\sgn (\sigma \tau \sigma^{-1}) = \sgn(\sigma)\sgn(\tau)\sgn(\sigma^{-1}) = \sgn(\sigma)^2\sgn(\tau) = 1^2\cdot 1 = 1$ 
    in zato
    $\sigma A_n \sigma^{-1} = A_n$.

    Očitno je $A_n \cap C_2 = \{e\}$, saj je $C_2$ ciklična grupa izomorfna eni transpoziciji, ki je liha permutacija.
    Dodatno pa vidimo $A_nC_2 = S_n$ saj na trivialen način dobimo vse sode permutacije, 
    za vse lihe permutacije $\sigma \in S_n$ pa velja kot $\tau \sigma = a \in A_n$, 
    torej $\sigma = \tau^{-1} a \in A_nC_2$.
\end{dokaz}

\begin{zgled}
    
\end{zgled}


\subsection{Orbite in delovanje grup}


\begin{definicija}
    Grupa $G$ je primitivna, če deluje tranzitivno na množici $X$ in če 
    delovanje grupe $G$ na množici $X$ ne ohranja nobene netrivialne particije množice $X$.
\end{definicija}





%%%%%%%%%%%%%%%%%%%%%%%%%%%%%%%%%%%
\begin{definicija}
    Levi koset podgrupe $H$ v grupi $G$ je množica 
    $gH = \{gh \mid h \in H\}$, kjer je $g$ poljuben element iz $G$. 
\end{definicija}
\begin{trditev}
    Levi koseti tvorijo disjunktno dekompozicijo grupe $G$ na podmnožice, ki so vse enake velikosti
    kot $H$.
\end{trditev}
\begin{dokaz}
    Ker je $H$ podgrupa $G$, vsebuje enoto $e$. Posledično je 
    $\bigcup_{g \in G} gH \supseteq \bigcup_{g \in G} ge = \bigcup_{g \in G} g = G$.

    Dalje $H$ in $gH$ imata enako število elementov, saj je $h \mapsto gh$ bijekcija med njima in 
    $x \mapsto g^{-1} x$ njen inverz. 

    Pokažimo še disjunktnost kosetov. Naj bosta $g_1 H$ in $g_2 H$ dva različna leva koseta, 
    katerih presek vsebuje $x$ potem je $x = g_1 h_1$ in hkrati $g_2 h_2$. Velja
    $g_1 = g_2 h_2 h_1^{-1}$, za poljuben $h$ je tako $g_1 h = g_2 h_2 h_1^{-1} h \in g_2 H$,
    kar  pomeni, da je $g_1 H \subseteq g_2 H$, simetrično dobimo tudi $g_2 H \subseteq g_1 H$, 
    torej $g_1 H = g_2 H$.
\end{dokaz}


\begin{definicija}
    Indeks podgrupe $H$ v grupi $G$ je število levih
    kosetov $H$-ja v $G$. Indeks podgrupe $H$ v grupi $G$ se označuje z $[G : H]$.
    Ker levi koseti tvorijo disjunktno dekompozicijo grupe $G$ na podmnožice, ki so si 
    enake v velikosti, za indeks velja $[G : H] = |G| / |H|$.
\end{definicija}

\begin{zgled}
Naj bo $G = \mathbb{Z}$ grupa celih števil z operacijo seštevanja. Naj bo $H = 2\mathbb{Z}$ 
podgrupa sodih celih števil. Potem ima $H$ dva koseta v $G$, 
namreč za vsak $g \in G$ je $g + H = \begin{cases}
\mathbb{Z} & \text{če je $g$ sodo} \\
\mathbb{Z} + 1 & \text{sicer}
\end{cases}$. Indeks podgrupe $H$ v grupi $G$ je torej $[G : H] = 2$.
\end{zgled}

\begin{izrek}[Lagrangev izrek]\label{iz:index_lagrange}
    Naj bo $H$ podgrupa $G$ in $K$ podgrupa $H$. Potem je 
    red grupe $G$ deljiv z redom grupe $H$ in $[G : K] = [G : H][H : K]$.
\end{izrek}
\begin{dokaz}
    Prva trditev sledi iz definicije indeksa, drugo pa pokažimo takole.
    Naj bo $S$ set predstavnikov kosetov za $K$ v $H$ in naj bo $T$ set predstavnikov
    kosetov za $H$ v $G$. Potem je $G = \bigsqcup_{t \in T} t H = \bigsqcup_{t \in T} \bigsqcup_{s \in S} tsK$.
    Posledično je set predstavnikov kosetov za $K$ v $G$ enak $TS = \{ts \mid t \in T, s \in S\}$, 
    kar pomeni, da je $[G : K] = |T||S| = [G : H][H : K]$.
\end{dokaz}

%%%%%%%%%%%%%%%%%%%%
\begin{definicija}
    Orbita elementa $x \in X$ je podmnožica elementov iz $X$, ki jih lahko dosežemo
    z delovanjem elementov iz grupe $G$. Pišemo $G \cdot x = \{g \cdot x \mid g \in G\}$.

    Množico vseh orbit delovanja grupe $G$ na množici $X$ imenujemo 
    kvocientna množica $X/G$. Označimo tudi $X/G_x = \{gG_x \mid g \in G\}$.
    Če ima grupa eno samo orbito, pravimo, da grupa deluje tranzitivno na množici $X$.
\end{definicija}

\begin{definicija}
    Naj bo $G$ grupa, in $g$ njen element. Pravimo, da $x \in X$ fiksna točka elementa $g$,
    če je $g \cdot x = x$. Množico vseh fiksnih točk elementov iz $G$, ki fiksirajo element $x$,
    imenujemo stabilizatorska podgrupa $G$ elementa $x$, pišemo 
    $G_x = \{g \in G \mid g \cdot x = x\}$.
\end{definicija}

\begin{trditev}\label{tr:stabilizator_sosedstvo}
    Naj bo $H$ graf in $V(H)$ množica njegovih vozlišč. Naj bo $Q \subseteq V(H)$
    podmnožica vozlišč. Stabilizatorska grupa $G_Q$ je podgrupa avtomorfizmov grafa $H$ in
    vsak njen element $g$ ohranja sosedstvo množice $Q$, 
    torej $g(N(Q)) = N(Q)$.
\end{trditev}
\begin{dokaz}
    Vzemimo $v \sim Q$. Ker je $g$ avtomorfizem, je 
    $g(v) \sim g(Q)$. Ker dalje $g$ fiksira množico $Q$, 
    je $g(Q) = Q$, in zato $g(v) \in N(Q)$, kar pomeni, da $g(N(Q)) \subseteq N(Q)$.
    Ker je $g$ bijekcija, je $g^{-1}$ tudi avtomorfizem, ki fiksira množico $Q$,
    torej $g^{-1}(N(Q)) \subseteq N(Q)$, kar pomeni, da $N(Q) \subseteq g(N(Q))$.
\end{dokaz}

\begin{izrek}[Izrek o orbitah in stabilizatorjih]\label{iz:orbite_stabilizator}
    Naj bo $x \in X$ fiksen in naj bo $f: G \to X$ podana z $g \mapsto g \cdot x$. 
    Potem je $f(G) = G \cdot x$ orbita elementa $x$ in 
    $f^{-1}(x) = G_x$ stabilizatorska podgrupa elementa $x$.
    Če imata dva elementa enako sliko, potem velja
    \begin{align*}
    f(g) = f(h) &\iff g\cdot x = h\cdot x \iff g^{-1}h\cdot x = x \iff g^{-1}h \in G_x \iff h \in gG_x 
    \end{align*}
    
    Torej, če sta $g$ in $h$ elementa enakega koseta (tukaj $gG_x$) 
    za stabilizatorsko grupo $G_x$, potem imata enako sliko, in obratno.
\end{izrek}

\begin{izrek}\label{iz:orbita_stabilizator_red}
    Če je $G$ končna potem red orbite elementa $x$ enak številu kosetov stabilizatorske
    podgrupe $[G : G_x]$, kar je enako $|G| / |G_x|$.
\end{izrek}
\begin{dokaz}
    Izrek sledi neposredno iz izreka o orbitah in stabilizatorjih~\ref{iz:orbite_stabilizator}
    in iz lagrangevega izreka o številu kosetov~\ref{iz:index_lagrange}.
\end{dokaz}


\begin{trditev}
    Orbite delovanja grupe $G$ na množici $X$ tvorijo particijo množice $X$.
\end{trditev}
\begin{dokaz}
    Pokažemo lahko, da so orbite ekvivalenčni razredi neke ekvivalenčne relacije,
    po tem bomo vedeli, da nujno razdelijo $X$ na disjunktne podmnožice, katerih unija je $X$.

    Definiramo relacijo $\sim$ na $X$ z 
    $x \sim y \iff \exists g \in G: g \cdot x = y \iff y \in G\cdot x$.
    Taka relacija je refleksivna $(x \sim x)$, saj je za enoto $e\in G$ $e \cdot x = x$. 
    Je simetrična $(x\sim y \implies y \sim x)$, saj če velja $x\sim y$,
    potem obstaja $g\in G$, tak da $g\cdot x = y$. Dobimo $x = g^{-1} \cdot y$.
    In je tranzitivna $(x\sim z \land z\sim y \implies x \sim y)$, saj če 
    velja $x\sim y$ (torej $y = g\cdot x$) in $z\sim y$ (torej $y = h \cdot z$), 
    potem je $h^{-1}g \cdot x = z$, kar pomeni, da je $x \sim z$.
\end{dokaz}




\begin{zgled}\label{zg:petersen_avtomorfizmi}
    Teorijo orbit in stabilizatorjev lahko uporabimo za štetje avtomorfizmov Petersenovega
    grafa. Petersenov graf je vozliščno tranzitiven, kar pomeni, da lahko s simetrijo 
    katerokoli vozlišče pretvorimo v katerokoli drugo vozlišče. Ker ima Petersenov 
    graf $10$ vozlišč, je orbita poljubnega vozlišča $v$
    reda $|G\cdot v| = 10$. Zdaj pa si zamislimo, da $v$ fiksiramo na mestu, 
    ker je Petersenov graf $3$-regularen, ima $v$ natanko $3$ sosede, ki jih lahko vsak 
    avtomorfizem, ki fiksira $v$ zgolj premeša. Premeša jih lahko na $3!$ načinov, kar 
    pomeni, da je stabilizatorska grupa $G_v$ reda $6$. 
    Po izreku o orbitah in stabilizatorjih~\ref{iz:orbite_stabilizator} je 
    red avtomorfizmov Petersenovega grafa enak $|G| = |G\cdot v||G_v| = 10 \cdot 6 = 60$.
\end{zgled}
\begin{zgled}
    Predstavljajmo si graf, ki je sestavljen iz dveh trikotnikov, ki sta povezana z eno samo
    povezavo, recimo mu $H$. Vozlišča razdelimo na dve začetni orbiti $O_1$ in $O_2$, 
    glede na njihovo stopnjo.
    Tako vzamemo $O_1 = \{v_1, w_1\}$ vozlišči na mostu in $Q_2 = \{v_2, v_3, w_2, w_3\}$
    vozlišča stopnje $2$.

    Izberimo zaporedje vozlišč $v_1, v_2, w_2$ in sestavimo verigo podgrup, kjer vsaka fiksira
    eno točko več:
    \[\aut(H) = G \geq G^{(1)} \geq G^{(2)} \geq G^{(3)} = \{e\}\]
    kjer je $e$ enota, torej avtomorfizem $\id$.
    
    Po Lagrangevem izreku~\ref{iz:index_lagrange} je 
    $|G| = [G : G^{(1)}][G^{(1)} : G^{(2)}][G^{(2)} : G^{(3)}]$. 

    V celotnem grafu se vozlišče $v_1$ lahko preslika le v vozlišče, ki je enake stopnje,
    torej zgolj samo vase ali v $w_1$. Ker sta oba strukturno enaka, je $[G : G^{(1)}] = 2$.

    Ko fiksiramo $v_1$, postane celoten desni del grafa (trikotnik $w_1, w_2, w_3$) 
    nepremičen glede na levo stran, saj mostu ne moremo več zrcaliti.
    V levem trikotniku pa se nahajata še $v_2$ in $v_3$, ki ju lahko zamenjamo z zrcaljenjem
    levega trikotnika. Tako je $[G^{(1)} : G^{(2)}] = 2$.

    Enako kot smo storili za $v_2$ in $v_3$ lahko zdaj storimo še za $w_2$ in $w_3$, 
    kar pomeni, da je $[G^{(2)} : G^{(3)}] = 2$, in posledično 
    $|\aut(H)| = 2 \cdot 2 \cdot 2 = 8$.
\end{zgled}

\section{Mooreovi grafi}
% definicija, slike
Mooreovi grafi so ekstremalni grafi, ki glede na dan premer ter dano 
maksimalno stopnjo dosežejo maksimalno število vozlišč. Leta 1960 sta jih tako 
poimenovala Hoffman in Singelton po Edwardu F. Mooru, ki je prvi zastavil vprašanje o 
opisu in klasifikaciji teh struktur.

\begin{definicija}\label{def:moore}
Naj bo $G$ graf s premerom $D$ ter maksimalno stopnjo $\Delta$.
Pravimo, da je $G$ Mooreov graf, če njegovo število vozlišč doseže zgornjo teoretično mejo:
\begin{equation}\label{eq:moorebound}
    |V| = 1 + \Delta + \Delta (\Delta - 1) + \ldots + \Delta (\Delta - 1)^{D - 1}.
\end{equation}
V nadaljevanju bomo simbola $D$ in $\Delta$ vselej uporabljali za označevanje parametrov 
Mooreovega grafa.
\end{definicija}

\begin{opomba}
Vrednost v enačbi~\ref{eq:moorebound} imenujemo Mooreova meja. Njen pomen postane jasen, 
če opazujemo strukturo grafa iz izbranega vozlišča $v \in V$:
na razdalji $1$ od $v$ je največ $\Delta$ sosedov, na vsaki naslednji razdalji $i$ 
(kjer $1 < i < D)$ pa ima vsako vozlišče na voljo največ $\Delta - 1$ sosedov,
ki niso bili obiskani v prejšnjih korakih.
Mooreov graf je torej tisti mejni primer, kjer se veje sosedstva nikoli ne prekrivajo 
dokler ne dosežemo premera $D$. 
\end{opomba}

\subsection{Lastnosti Mooreovih grafov}

\begin{trditev}\label{tr:mooreregularpaths}
Mooreov graf je regularen stopnje $\Delta$, njegova ožina je enaka ${2D+1}$,
ter za poljubni dve vozlišči $u$ in $v$
velja, da med njima obstaja natanko ena najkrajša pot (dolžine največ~$D$).
\end{trditev}

\begin{figure}[ht]\label{fig:mooretree}
    \centering
    \def\svgwidth{\columnwidth}
    \resizebox{0.9\textwidth}{!}{%
        \import{../../assets/}{moore-bound.tex}
    }
    \caption{Lokalna struktura Moorejevega grafa.}
\end{figure}


\begin{dokaz}
V Mooreevem grafu izberimo poljubno vozlišče $v$, če je stopnja tega vozlišča 
manjša od $\Delta$, potem Mooreova meja ne more biti dosežena, saj velja: 
\[|V(G)| \leq 1 + \deg(u)(\Delta - 1) + \ldots + \deg(u) (\Delta - 1)^{D - 1}.\]


Hkrati se o enoličnosti najkrajše poti prepričamo tako, da iz poljubnega vozlišča $v$
tvorimo inducirano drevo (kakor je prikazano na sliki~\ref{fig:mooretree}). To drevo je
očitno tudi drevo najkrajših poti (in sicer enolično).

Neposredno iz konstrukcije Mooreeve meje (glej~\ref{fig:mooretree}) je razvidno tudi,
da je ožina Mooreovega grafa enaka $g=2D+1$. Krajših ciklov v grafu ni.
\end{dokaz}


V nadaljevanju bomo zasledili, da so Mooreovi grafi deležni številnih dodatnih
omejitev, kar bo utemeljilo tudi njihovo redkost, namreč Mooreovi grafi obstajajo
zgolj za omejene nabore parametrov $D$ in $\Delta$.

\subsubsection{Matrična enačba Mooreovih grafov}

%citiraj izpeljavo iz uvoda v teorijo
Naj bo $A$ sosednostna matrika Mooreovega grafa stopnje $\Delta$. 
Definirajmo matriko $G^{\Delta, i}$ 
tako, da njeni elementi $G^{\Delta, i}_{uv}$ podajo število poti  
dolžine največ $i$ med $u$ in $v$. Ker so edine poti dolžine $0$ vozlišča sama, 
velja $G^{\Delta, 0} = \II$, za poti dolžine največ $1$ pa $G^{\Delta, 0} = A + \II$.
Produkt $AG^{\Delta, i-1}$ prešteva sprehode dolžine največ $i$, ki jih dobimo tako, 
da potem dolžine $(i-1)$ dodamo še en korak. Pri tem upoštevamo, da se nekateri
sprehodi iz $(i-1)$~-tega nivoja vrnejo nazaj v nivo $(i-2)$.
Ker je Mooreov graf $\Delta$~-regularen (\ref{tr:mooreregularpaths}), za vsak par vozlišč $(u, v)$ obstaja natanko
$(\Delta - 1)G^{\Delta, i-2}_{uv}$ takšnih sprehodov, ki niso poti in jih je treba 
odšteti. S tem razmislekom smo ravnokar dokazali naslednjo trditev.

\begin{trditev}\label{tr:moorematrix}
Sosednostna matrika Mooreovega grafa zadošča matrični enačbi 
\[G^{\Delta, D} = J,\]
kjer je 
$J$ matrika samih enic, $G^{\Delta, D}$ pa dobljena rekurzivno z:
\begin{equation}
\begin{cases} 
G^{\Delta, 0} = I \\
G^{\Delta, 1} = A + I \\
G^{\Delta, i} = A G^{\Delta, i-1} - (\Delta-1) G^{\Delta, i-2} & \text{za } 2 \leq i \leq D
\end{cases}
\end{equation}
\end{trditev}

\begin{opomba}
    Velja celo to, da regularni graf z defektom $\delta$ zadošča enačbi
    $G^{\Delta, D} = J + B$, kjer je $B$ matrika z 
    stolpičnimi in vrstičnimi seštevki enaki $\delta$ \cite[str. 15]{MillerSiran2013}.
    Več o grafih z defektom bomo povedali kasneje~(\ref{sec:moorerazsiritev}).
\end{opomba}

% Tu lahko rečeš da dodatno ta enačba celo posplošuje na grafe z defektom delta.

\subsubsection{Izrek Hoffman-Singeltona}
Naj bo premer $D$ enak $2$, ter stopnja $\Delta$ večja od $1$. Tedaj za matriko 
sosednosti Mooreovega grafa (po trditvi \ref{tr:moorematrix}) velja:
\begin{equation}\label{eq:moorematrix2}
A^2 + A - (\Delta-1)\II = J,
\end{equation}
kjer je $\II$ identiteta in $J$ matrika samih enic.
Iz definicije~\ref{def:moore} je razvidno, da je red Mooreovega grafa enak
$\Delta^2 + 1$. Velja tudi, da je Mooreov graf premera $2$ krepko regularen 
s parametri $(\Delta^2 + 1, \Delta, 0, 1)$, kar pomeni, da je vsak par sosednjih vozlišč 
povezan z $0$ skupnimi sosedi (glej omejitev ožine~\ref{tr:mooreregularpaths}), 
vsak par nesosednjih vozlišč pa je povezan z $1$ skupnim sosedom (glej omejitev premera~\ref{def:moore}).

Matrika $J$ je dimenzije $(\Delta^2 + 1) \times (\Delta^2 + 1)$.
Ker velja $J\cdot\mathbf{1} = (\Delta^2 + 1)\cdot\mathbf{1}$, je $(\Delta^2 + 1)$ njena 
lastna vrednost in $\mathbf{v}_1 = \mathbf{1}$ njen lastni vektor. 
Ker je $\rank(J) = 1$, in s tem $\dim(\ker(J)) = \Delta^2$ \cite{enwiki1344320194}, je njena večkratnost enaka $1$, ostale
lastne vrednosti pa so tako enake $0$, z večkratnostjo $\Delta^2$,
za lastne vektorje 
$(\mathbf{v}_2, \ldots, \mathbf{v}_{\Delta^2}) = (\mathbf{e}_1 - \mathbf{e}_2, \ldots, \mathbf{e}_1 - \mathbf{e}_{\Delta^2 + 1})$,
kjer je $\mathbf{e}_i$ $i$-ti standardni bazni vektor.

Z množenjem matrične enačbe~\ref{eq:moorematrix2} z lastnimi vektorji $\mathbf{v}_i$ dobimo naslednji dve enačbi:

\begin{align}
x^2 + x - (\Delta - 1) = \Delta^2 + 1 \label{eq:karakteri1}\\
x^2 + x - (\Delta - 1) = 0 \label{eq:karakteri2}
\end{align}


Iz enačbe~\ref{eq:karakteri1} dobimo lastno vrednost $\Delta$ z večkratnostjo $1$,
Iz enačbe~\ref{eq:karakteri2} pa dobimo dve lastni 
vrednosti $r$ in $s$ enaki $(-1 \pm \sqrt{4\Delta - 3})/2$, s pripadajočima večkratnostima $a$ in $b$, 
kjer $1 + a + b = 1 + \Delta^2$.
Ker je sled matrike sosednosti vedno enaka $0$, in ker velja 
$\sled(A) = \sum_{\lambda \text{ je lastna vr.}} \lambda$,
je $\Delta + a r + b s = 0$.

Če sta $r$ in $s$ iracionalni števili je $a = b$, potem 
$\Delta + a r + b s = 0 \implies \Delta - \Delta^2 / 2 = 0 \implies \Delta = 2$.
Sicer pa sta $r$ in $s$ celi števili, namreč $4\Delta - 3$ je vedno liho število, posledično 
je tudi njegov koren vedno lih. Tako nastavimo $4\Delta - 3 = m^2$ za nek $m \in \ZZ_+$, 
oziroma $\Delta = \frac{m^2 + 3}{4}$, od tod dobimo: 
$\Delta + a r + b s = 0 \implies m^5 + m^4 + 6m^3 - 2m^2 + 9m - 15 = 0$.
Ker ima slednji moničen polinom celoštevilske koeficiente, so po izreku o 
celoštevilskih ničlah~\cite{enwiki1346919033}, njegove rešitve deljitelji števila $15$,
torej $m \in \{1, 3, 5, 15\}$, kar nam da možne rešitve za $\Delta$ enake $\{1, 3, 7, 57\}$.

Iz zgornjih dveh enačb, ki sta linearni v $a$ in $b$, izvemo tudi, da sta $a$ in $b$ enaka 
$\frac{1}{2} [\Delta^2 \mp \frac{2\Delta - \Delta^2}{\sqrt{4\Delta - 3}}]$.

\begin{izrek}[Hoffman-Singeltonov izrek]\label{iz:hoffmansingeltonizrek}
    Mooreovi grafi z premerom $D = 2$ in stopnjo $\Delta > 1$ obstajajo le za $\Delta \in \{2, 3, 7, 57\}$.
    Njihove lastne vrednosti so $\Delta$ z večkratnostjo $1$, 
    $\frac{1}{2}(-1 + \sqrt{4\Delta - 3})$ z večkratnostjo 
    $\frac{1}{2} [\Delta^2 - \frac{2\Delta - \Delta^2}{\sqrt{4\Delta - 3}}]$, ter 
    $\frac{1}{2}(-1 - \sqrt{4\Delta - 3})$ z večkratnostjo
    $\frac{1}{2} [\Delta^2 + \frac{2\Delta - \Delta^2}{\sqrt{4\Delta - 3}}]$.
\end{izrek}

\begin{opomba}
Do matrične enačbe, ter spektra grafa, bi lahko prišli tudi z uporabo dejstva, 
da so Mooreovi grafi premera $2$ krepko regularni s parametri
$(n, k, \lambda, \mu) = (\Delta^2 + 1, \Delta, 0, 1)$. Namreč
krepko regularni grafi zadoščajo matrični enačbi $A^2 = k\II + \lambda A + \mu(J - \II - A)$,
ter spekter matrike $A$ je znan~\cite{enwiki1342817669}.
\end{opomba}


\subsection{Obstoj Mooreovih grafov}

Hitro lahko sestavimo nekaj trivialnih Mooreovih grafov, namreč, če je premer $D=1$, potem 
je Mooreov graf poln graf $K_{\Delta + 1}$ za vsako stopnjo $\Delta$, če je stopnja 
$\Delta = 1$, potem je Mooreov graf ciklični graf $C_{2D + 1}$ za vsak premer $D$.

Hoffman-Singeltonov izrek~\ref{iz:hoffmansingeltonizrek} pa nam pove, da za premer $D=2$ obstajajo 
največ štirje Mooreovi grafi, in sicer za stopnje $\Delta \in \{2, 3, 7, 57\}$.
Izkaže se, da so prvi trije Mooreovi grafi unikatni, imenovani Ciklični graf $C_5$,
Petersenov graf, ter Hoffman-Singeltonov graf, medtem ko obstoj grafa valence $57$ ostaja odprt problem.

Predstavimo sedaj netrivialne primere podrobneje, v naslednjih nekaj sekcijah.

\subsubsection{Petersenov graf}

Konstrukcija Petersenovega grafa je preprosta, in bi lahko bila puščena bralcu kot vaja, vendar jo bomo vseeno 
predstavili, in sicer na več različnih načinov, saj nam bo služila kot dobra vaja za razumevanje 
kasneje kompleksnejših konstrukcij.

Začnimo s konstrukcijo, ki izhaja iz drevesnega grafa~(glej \ref{def:moore}). 
Vzemimo neko začetno vozlišče $v$, ter ga povežimo z $\Delta=3$ sosedi, ki jih poimenujemo $v_1, v_2, v_3$.
Vsakega izmed teh povežimo še z dvema novima vozliščema, ki jih poimenujemo 
$v_{11}, v_{12}, v_{21}, v_{22}, v_{31}, v_{32}$. V tem trenutku v rokah držimo drevesni graf. Sedaj moramo
njegove liste povezati tako, da ne ustvarimo ciklov dolžine $3$ ali $4$~(glej \ref{tr:mooreregularpaths}). 
Brez škode za splošnost lahko 
povežemo $v_{11}$ z $v_{21}$ ter z $v_{31}$, $v_{12}$, po tem smo prisiljeni povezati $v_{12}$ z 
$v_{22}$ in $v_{32}$, brez drugih izbir tako ostanemo do konca konstrukcije, povežemo $v_{21}$ z
$v_{32}$, ter $v_{22}$ z $v_{31}$.

\begin{figure}[ht]\label{fig:petersendrevo}
    \centering
    \def\svgwidth{\columnwidth}
    \resizebox{0.8\textwidth}{!}{%
        \import{../../assets/}{petersen-drevo.tex}
    }
    \caption{Drevesna konstrukcija Petersenovega grafa.}
\end{figure}

Alternativno lahko uporabimo dejstvo, da je ožina takega grafa enaka $5$~(glej \ref{def:moore}), torej v takem grafu obstaja
cikel $C_5$. Začnimo zato z ciklom $C_5$. Ker vsakemu vozlišču v tem ciklu manjka še en sosed, mu ga dodamo.
Tu opazimo, da vozlišča cikla $C_5$ ne smejo imeti skupnih sosedov (izven tega cikla). Če se sedaj osredotočimo
na eno izmed novo dodanih vozlišč, vidimo, da je lahko povezano le z dvema vozliščema, ki sta soseda nasprotnih dveh 
vozlišč v ciklu $C_5$. Hkrati mora biti povezano z natanko dvema. Povežemo ga z njima. Simetrično nato storimo 
za druge. Taki konstrukciji
rečemo konstrukcija iz petkotnika, poslužili se je bomo tudi pri konstrukciji Hoffman-Singeltonovega grafa.

\begin{figure}[ht]\label{fig:petersenpetkotnik}
    \centering
    \def\svgwidth{\columnwidth}
    \resizebox{0.8\textwidth}{!}{%
        \import{../../assets/}{petersen-petkotnik.tex}
    }
    \caption{Konstrukcija Petersenovega grafa iz petkotnika.}
\end{figure}

Petersenov graf lahko konstruiramo tudi s pomočjo simetrične grupe \newline $S_5 = \{1, 2, 3, 4, 5\}$, 
kjer so vozlišča grafa predstavljena z $2$-elementnimi podmnožicami množice $S_5$, medtem ko so povezave
med vozlišči predstavljene z njihovo disjunktnostjo. Pišemo lahko 
$V(\text{Petersenov graf}) = \{\{a, b\}\mid\{a, b\}\subset S_5\}$ ter $uv \in E(\text{Petersenov graf}) \iff u \cap v = \emptyset$. 
Ta konstrukcija res poda Mooreov graf na $10$ vozliščih, stopnje $3$ in premera $2$, 
saj ima ustvarjen graf $\binom{5}{2} = 10$ vozlišč, 
za poljubno vozlišče $\{a, b\}$ obstajajo natanko $\binom{3}{2} = 3$ pari, ki so z njim disjunktni, in zato sosednji,
ter ker so cikli krajši od $5$, po konstrukciji, nemogoči.
Iz te konstrukcije je razvidno, da je Petersenov graf Kneserjev graf $\kneser(5, 2)$, torej tudi, da je 
vozliščno tranzitiven, in da je
$\aut(\text{Petersenov graf}) \cong S_5$.



\subsubsection{Hoffman-Singeltonov graf}

Naj bo $H$ graf, ki ga iščemo, torej Mooreov graf premera $2$ ter stopnje $7$.
Tak graf sestoji iz $7\cdot 7 + 1 = 50$ vozlišč (\ref{def:moore}) in ima 
ožino $5$ (\ref{tr:mooreregularpaths}).
Po korakih Conga Fana in Allena J. Schwenka \cite{fan1993structure} dokažimo,
da tak graf, ki mu pravimo Hoffma-Singeltonov graf, obstaja in je enoličen.

Ker ima graf ožino $5$ v njem najdemo vsaj en cikel $C_5$. Označimo ga s
$Q=v_1 v_2 v_3 v_4 v_5 v_1$ in mu recimo centralni cikel. 
V nadaljevanju bomo opazovali njegove sosede
in s tem gradili strukturo grafa $H$.

\begin{trditev}\label{tr:gruce_preparation}
    Poljubno vozlišče centralnega cikla $v \in Q$ ima natanko $7 - 2 = 5$ sosedov nevsebovanih v $Q$.
\end{trditev}
\begin{dokaz}
    Če to ne bi veljalo, bi med vozlišči grafa $Q$ našli cikel $C_3$, kar bi 
    bilo v nasprotju s predpostavko o ožini.
\end{dokaz}

Sedaj si te sosede oglejmo podrobneje. 
Rečemo, da vozlišča grafa $H$, ki so sosedna vozlišču $v_i \in Q$, vendar ne ležijo v $Q$,
pripadajo gruči $\gr_i$. Dodatno bomo rekli, da je $v_i$ center svoje gruče $\gr_i$. 
Matematično bi to definirali kot:
\begin{equation*}
\forall v_i \in Q: \; \gr_i = \{\, u \in H \setminus Q \mid u \sim v_i\,\} = N(v_i) \setminus V(Q).
\end{equation*}
Pri taki definiciji je dobro uvideti, da unija takih gruč pokrije ravno vse sosede vozlišč v $Q$,
in da je moč vsake take gruče po trditvi~\ref{tr:gruce_preparation} enaka 5.

\begin{figure}[ht]
    \centering
    \def\svgwidth{0.7\columnwidth}
    \import{../../assets/}{hoffman-singelton-proof.pdf_tex}
    \caption{Vizualni prikaz strukture gruč v grafu $H$.}
\end{figure}

\begin{trditev}\label{tr:lastnosti_gruc}
Gruče, kot definirane zgoraj, v grafu $H$ zadoščajo naslednjim lastnostim:
\begin{enumerate}[label=(\roman*)]
\item Gruče so si paroma disjunkne: \(\forall i, j: \;\gr_i \cap \gr_j = \emptyset\).
    \label{itm:g1}
\item Vozlišča znotraj posamezne gruče so nepovezana: 
\(\forall i. \forall u, v \in \gr_i: \; u \not\sim v\).  
    \label{itm:g2}
\item Če sta centra \(v_i\) in \(v_j\) sosedna v \(Q\), med gručama ni povezave,
sicer pa je vsako vozlišče gruče \(\gr_i\) sosedno največ enemu vozlišču 
gruče \(\gr_j\). Oziroma: \\ \(\forall i, j. \forall u\in \gr_i:\; \left( (v_i \sim v_j 
\implies \gr_i \not\sim \gr_j) \lor  (|N(u) \cap \gr_j| \leq 1)\right)\)
    \label{itm:g3}
\end{enumerate}
\end{trditev}

\begin{dokaz}
Vzemimo vozlišči $u, w \in \gr_i$, ter vozlišče $z \in \gr_j$, kjer $i \neq j$.
Pri dokazih uporabimo metodo protislovja, kjer se dokaz zaključi ko najdemo 
vsaj en cikel dolžine manj kot 5. To bi namreč bilo v nasprotju z ožino grafa.
\begin{enumerate}[label=(\roman*)]
\item Predpostavimo, da obstajata dve gruči, katerih presek ni prazen.
Brez škode za splošnost predpostavimo $u = z \in \gr_i \cap \gr_j$.
Tedaj v grafu $H$ obstaja cikel $u v_i \ldots v_j u$, ki je bodisi $C_3$ ali $C_4$.
\item Naj obstaja gruča $\gr_i$ ter znotraj nje vozlišči $u$ ter $w$, ki sta povezani.
Tedaj v grafu obstaja cikel $C_3 = u v_i w u$.
\item Naj bosta najprej centra gruč $v_i$ ter $v_j$ sosedna, ter $u$ in $z$
povezana. Tedaj v grafu $H$ obstaja cikel $C_4 = u v_i v_j z u$.
Predpostavimo sedaj, da $v_i \not \sim v_j$, ter da je vozlišče $z$ povezano 
tako z $u$ kot z $w$. Tedaj obstaja cikel $C_4 = z u v_i w z$.
\end{enumerate}
\end{dokaz}

\begin{posledica}\label{po:enolicna_povezava_iz_nq_v_q}
    Poljubno vozlišče $u$ iz $N(Q)$ je povezano z natanko enim vozliščem iz centralnega
    cikla $Q$. Namreč po definiciji gruč je unija vseh njih ravno $N(Q)$, 
    po trditvi~\ref{tr:lastnosti_gruc}\ref{itm:g1} pa, ker so si te gruče disjunktne,
    velja da je taka unija disjunktna.
    Torej je $u$ del natanko ene gruče, ki je po definiciji povezana z natanko enim vozliščem iz $Q$.
\end{posledica}

\begin{posledica}\label{po:lastnosti_gruc}
Izvedeli smo tudi, da je maksimalna stopnja
podgrafa $H[N(Q)]$ induciranega z množico sosedov $Q$, označeno z $N(Q)$, enaka $2$. 
Po točki~\ref{itm:g3} trditve~\ref{tr:lastnosti_gruc} namreč velja, da je vsako vozlišče dane 
gruče $\gr_i$ lahko povezano zgolj z gručami katerih centri si niso sosedni z $v_i$.
Taki gruči sta dve, ter do vsake izmed njiju lahko vodi največ ena povezava.
\end{posledica}

Ogledali smo si strukturo centralnega cikla $Q$ ter strukturo njegovih sosedov $N(Q)$,
katerih je, mimogrede, natanko $5\cdot 5 = 25$ 
(glej \ref{tr:gruce_preparation} in \ref{tr:lastnosti_gruc}\ref{itm:g1}).

Da bomo o strukturi grafa $H$ lahko povedali še kaj več, si oglejmo še preostala vozlišča, 
katerih je $50 - 5 - 25 = 20$. Ker so le ta, na razdalji natanko $2$ od vozlišč centralnega cikla $Q$,
jih prikladno označimo z $N_2(Q)$.

\begin{trditev}\label{tr:st_povezav_iz_n2q}
    Naj bo $q$ vozlišče iz množice $N_2(Q)$, tedaj ima lahko z vozlišči iz $N(Q)$ največ $5$ povezav.
\end{trditev}
\begin{dokaz}
    Če bi jih imelo več, bi to pomenilo, da je $q$ povezano z neko gručo $\gr_i$ dvakrat,
    recimo z $u$ in $w$. Potem bi v grafu $H$ obstajal cikel $C_4 = q u v_i w q$,
    kar bi bilo v nasprotju z ožino grafa.
\end{dokaz}

\begin{trditev}\label{tr:st_sosedov_v_n2q}
Vsako vozlišče $u$, ki pripada gruči $\gr_i$, ima natanko $4$ sosede v $N_2(Q)$. Poleg tega
vozlišča iz iste gruče nimajo skupnih sosedov v $N_2(Q)$. Formalno to zapišemo kot:
\[
\forall i:\; \forall u \neq w \in \gr_i:\ |N(u) \cap N_2(Q)| = 4 \ \land \ N(u) \cap N(w) \cap N_2(Q) = \emptyset,
\]
\end{trditev}

\begin{dokaz}
Kot v dokazu trditve~\ref{tr:lastnosti_gruc} bomo tudi tu vzeli vozlišči $u, w \in \gr_i$,
ter vozlišče $z\in\gr_j$, kjer $i \neq j$. Tudi tu bomo uporabili metodo protislovja.

Opazujmo najprej vozlišče $w$. Tako vozlišče ima vsaj 4 sosede v $N_2(Q)$, 
saj je njegova stopnja enaka $7$, povezan pa je z vozliščem 
$v_i \in Q$ ter po posledici~\ref{po:lastnosti_gruc} še z največ dvema
vozliščema iz $N(Q)$. Hkrati pa v $N_2(Q)$ nima več kot 4 sosede, kar pokažemo z enostavnejšim 
protislovjem. Predpostavimo lahko nasprotno, potem bi iz množice vozlišč $N(Q)$ 
v množico vozlišč $N_2(Q)$ vodilo vsaj 
$24\cdot 4 + 1\cdot 5 = 101$ povezav, kar je več kot $20\cdot 5 = 100$, ki bi vodilo v obratno 
smer (namreč vozlišče iz $N_2(Q)$ ima lahko največ $5$ sosedov v $N(Q)$ 
\ref{tr:st_povezav_iz_n2q}).
\end{dokaz}

\begin{posledica}\label{po:regularnosti_podgrafov}
    S tem smo pokazali tudi, da je podgraf $H[N(Q)]$ induciran z množico $N_2(Q)$ 2-regularen ter,
    da ima vsako vozlišče iz $N_2(Q)$ natanko 5 sosedov v $N(Q)$. Posledično je tudi $H[N_2(Q)]$
    2-regularen.
\end{posledica}

\begin{dokaz}
    Vozlišče iz podgrafa $H[N(Q)]$ je po trditvi~\ref{tr:st_sosedov_v_n2q} povezano 
    z natanko 4 vozlišči iz $N_2(Q)$, torej iz množice $N(Q)$ v množico $N_2(Q)$
    vodi natanko $25\cdot 4 = 100$ povezav. V dokazu ravno te trditve pa smo pokazali,
    da lahko v nasprotno smer vodi toliko povezav zgolj, če ima vsako vozlišče iz $N_2(Q)$ 
    natanko $5$ sosedov v $N(Q)$, od koder naravno sledi, da je $H[N_2(Q)]$ $2$-regularen.
    
    Poleg tega pa je, po 
    posledici~\ref{po:enolicna_povezava_iz_nq_v_q}, tako vozlišče iz $H[N(Q)]$ povezano z 
    natanko enim vozliščem iz centralnega cikla $Q$. Skupno ima tako vozlišče natanko $5$ povezav izven 
    $H[N(Q)]$ in zato točno $2$ povezavi znotraj $H[N(Q)]$ -- je $2$-regularen.

    Pri dokazu trditve~\ref{tr:st_sosedov_v_n2q} 
\end{dokaz}

\begin{trditev}\label{tr:unija_ciklov}
    Podgraf $H[N(Q)]$ induciran z množico $N(Q)$ je disjunktna unija ciklov, 
    katerih dolžina je večkratnik števila $5$.
\end{trditev}

\begin{dokaz}
    Tu se preprosto skličemo na posledico~\ref{po:regularnosti_podgrafov}, ki nam pove,
    da je $H[N(Q)]$ 2-regularen, kar je ekvivalentno temu, da je disjunktna unija ciklov.
    Lastnosti gruč~\ref{tr:lastnosti_gruc}\ref{itm:g3} pa nam dodatno povedo, da lahko vsak 
    cikel zaključimo šele v koraku, ki je večkratnik števila $5$. 
\end{dokaz}

Sedaj se osredotočimo na spektralno analizo grafa $H$ s pomočjo katere bomo 
lahko sklepali več o številu in velikosti ciklov, ki sestavljajo podgraf $H[N(Q)]$. 
Iz lastnosti Mooreovih 
grafov~\ref{iz:hoffmansingeltonizrek} vemo, da je njegov spekter 
enak $\{7^{(1)}, 2^{(28)}, (-3)^{(21)}\}$.

Pred nadaljevanjem je smiselno predstaviti lemo, ki govori o lastnih vrednostih cikličnih
grafov, ter splošno znan izrek, ki govori o lastnih vrednostih podgrafov.

\begin{lema}\label{le:cikellastne}
Naj bo dan cikel $C_n$ ter njegova sosednostna matrika $A$. Tedaj so lastne vrednosti 
matrike $A$ enake $2\cos(2\pi k/n)$, kjer $k \in \{0, 1, \ldots, n-1\}$. Posebno za 
vsako lastno vrednost velja $-2 \leq 2\cos(2\pi k/n) \leq 2$.
\end{lema}
\begin{dokaz}
    Naj bo $A$ sosednostna matrika cikla $C_n$. Potem velja:
    \begin{align*}
        A =  \begin{pmatrix} 
            0 & 1 & 0 & \dots & 1 \\ 
            1 & 0 & 1 & \dots & 0 \\ 
            0 & 1 & 0 & \dots & 0 \\ 
            \vdots & \vdots & \vdots & \ddots & \vdots \\ 
            1 & 0 & 0 & \dots & 0 \end{pmatrix}
        =  \begin{pmatrix} 
            0 & 0 & 0 & \dots & 1 \\ 
            1 & 0 & 0 & \dots & 0 \\ 
            0 & 1 & 0 & \dots & 0 \\ 
            \vdots & \vdots & \vdots & \ddots & \vdots \\ 
            0 & 0 & 0 & \dots & 0 \end{pmatrix} 
            +
            \begin{pmatrix} 
            0 & 1 & 0 & \dots & 0 \\ 
            0 & 0 & 1 & \dots & 0 \\ 
            0 & 0 & 0 & \dots & 0 \\ 
            \vdots & \vdots & \vdots & \ddots & \vdots \\ 
            1 & 0 & 0 & \dots & 0 \end{pmatrix}
        = P + P^{-1}.
    \end{align*}
    Kjer je $P$ permutacijska matrika, ki predstavlja rotacijo vozlišč 
    cikla $C_n$ v smeri urinega kazalca.
    Nadaljno velja:
    \begin{align*}
        \det(P - \lambda I) &= -\lambda \begin{vmatrix} 
            -\lambda & 0 & 0 & \dots & 0 \\ 
            1 & -\lambda & 0 & \dots & 0 \\ 
            0 & 1 & -\lambda & \dots & 0 \\ 
            \vdots & \vdots & \vdots & \ddots & \vdots \\ 
            0 & 0 & 0 & \dots & -\lambda  \end{vmatrix}
            + (-1)^{n+1} \begin{vmatrix} 
            1 & -\lambda & 0 & \dots & 0 \\ 
            0 & 1 & -\lambda & \dots & 0 \\ 
            0 & 0 & 1 & \dots & 0 \\ 
            \vdots & \vdots & \vdots & \ddots & \vdots \\ 
            0 & 0 & 0 & \dots & 1  \end{vmatrix} \\
            &= (-1)^n \lambda^n + (-1)^{n+1}
    \end{align*}
    Posledično je karakteristični polinom matrike $P$ enak $\lambda^n - 1$, katerega ničle 
    so $n$-ti koreni enote, torej so lastne vrednosti matrike $A$ enake 
    $\lambda_j = e^{i \frac{2\pi j}{n}} + e^{-i \frac{2\pi j}{n}} = 2\cosh(2\pi i j/n) = 2\cos(2\pi j/n)$.
\end{dokaz}


\begin{izrek}[Izrek o prepletanju lastnih vrednosti]\label{iz:interlacing}
Naj bo $G$ graf z $n$ vozlišči, ter $H$ njegov podgraf z $r$ vozlišči. Naj bo 
$A$ sosednostna matrika grafa $G$, ter $B$ sosednostna matrika grafa $H$. Naj bo
$\lambda_1 \geq \lambda_2 \geq \ldots \geq \lambda_n$ spekter matrike $A$, 
ter $\mu_1 \geq \mu_2 \geq \ldots \geq \mu_r$ spekter matrike $B$. Potem velja:
\[
\lambda_i \geq \mu_i \geq \lambda_{n - r + i} \quad \text{za vse } i \in \{1, 2, \ldots, r\}.
\]
\end{izrek}
\begin{dokaz}
    Dokaz izpustimo. Več lahko preberete v~\cite{enwiki1325771674}.
\end{dokaz}

Po lemi~\ref{le:cikellastne} vemo, da ima cikel $C_n$ natanko eno lastno vrednost enako $2$, 
in ker je $H[N(Q)]$ disjunktna unija ciklov, bo imel tak podgraf natanko toliko 
lastnih vrednosti enakih $2$, kolikor je ciklov v tej uniji.

Nato lahko izrek o prepletanju lastnih vrednosti~\ref{iz:interlacing} uporabimo
direktno na podgrafu $H[N(Q)]$, vendar dobimo prešibke rezultate.

Izkaže se, da je bolje uporabiti izrek o prepletanju lastnih vrednosti na 
podgrafu $H[v \cup N(Q)]$, ki ga dobimo, če iz grafa $H$ odstranimo vsa vozlišča, 
razen enega iz $Q$ in tistih iz $N(Q)$, nato pa opazovati spremembo multiplitete
lastne vrednosti $2$ ko iz podgrafa $H[v \cup N(Q)]$ odstranimo vozlišče $v$.

V ta namen navedimo naslednjo definicijo in lemo, ki nam bosta služili pri nadaljnji analizi.

\begin{definicija}
    Naj bo $H$ poljuben graf z množico vozlišč $V(H)$, za
    poljubno $W \subset V(H)$ neprazno, definiramo delni spoj, označen z $\oplus$, 
    kot $v \oplus (H, W)$, kjer je $v$ novo vozlišče, ki je povezano z vsemi 
    vozlišči iz $W$, ter z nobenim drugim vozliščem grafa $H$.
\end{definicija}


\begin{lema}\label{le:knjiga}
    Naj bo $G = \cup_{i=1}^k C_{n_i}$, ter $W \subset V(G)$ neprazna. Potem je 
    multipliteta lastne vrednosti $2$ v grafu $v \oplus (G, W)$ enaka 
    $k-1$.
\end{lema}
\begin{dokaz}
    Vredno je omeniti, da ima vsak cikel $C_n$ natanko eno lastno vrednost enako $2$ 
    (glej lemo~\ref{le:cikellastne}). Naj bodo $B_1, \ldots, B_k$ 
    sosednostne matrike ciklov $C_{n_1}, \ldots, C_{n_k}$
    ter $B$ sosednostna matrika grafa $G$.
    $B$ je bločno diagonalna matrika, katere bloki so 
    matrike $B_1, \ldots, B_k$. Posledično je multipliteta lastne vrednosti $2$ v 
    grafu $G$ enaka $k$.

    Naj bo dalje $A$ sosednostna matrika grafa $v \oplus (G, W)$,
    podrobneje pa nas zanima dimenzija lastnega prostora lastne vrednosti $2$; velja:
    \begin{align*}
        A - 2I &= \begin{bmatrix} 
            B_1 - 2I & 0 & 0 & \dots & x_{C_{n_1}} \\ 
            0 & B_2 - 2I & 0 & \dots & x_{C_{n_2}} \\ 
            0 & 0 & B_3 - 2I & \dots & x_{C_{n_3}} \\ 
            \vdots & \vdots & \vdots & \ddots & \vdots \\ 
            x_{C_{n_1}}^\top & x_{C_{n_2}}^\top & x_{C_{n_3}}^\top & \dots & -2  \end{bmatrix}
            = \begin{bmatrix}
            B - 2I & \mathbf{x} \\
            \mathbf{x}^\top & -2
            \end{bmatrix}
        \end{align*}
    Naj bo zdaj $\begin{bmatrix}
        \mathbf{v} \\
        t
    \end{bmatrix} \in \ker(A - 2I)$, potem je $0 = (A - 2I) \begin{bmatrix}
        \mathbf{v} \\
        t
    \end{bmatrix} = \begin{bmatrix}
        (B - 2I)\mathbf{v} + t \mathbf{x} \\
        \mathbf{x}^\top \mathbf{v} - 2t
    \end{bmatrix}$. Ker $(B-2I)\mathbf{v} + t \mathbf{x} = 0$, je
    $-t\mathbf{x} \in \im(B-2I)$. Ker je $(B-2I)$ simetrična je $\im(B-2I) = \ker(B-2I)^\perp$,
    torej za vsak $u \in \ker(B-2I)$ velja $-t\mathbf{x} \cdot u = 0$. 
    Ker je $\ker(B-2I) = \{\mathbf{1}_{C_{n_1}}, \ldots, \mathbf{1}_{C_{n_k}}\}$ 
    (\ref{le:cikellastne}), in ker $W \cap V(G) \neq \emptyset$, 
    obstaja $i \in \{1, 2, \ldots, k\}$, da za $\mathbf{x} \cdot \mathbf{1}_{C_{n_i}} > 0$
    in posledično iz 
    $-t\mathbf{x} \cdot \mathbf{1}_{C_{n_i}} = 0$ sledi $t=0$.

    Sistem se poenostavi v $\begin{bmatrix}
        (B-2I)\mathbf{v} \\
        \mathbf{x}^\top \mathbf{v}
    \end{bmatrix} = 0$. Od tod velja $\mathbf{v} \in \ker(B-2I)$, torej 
    $\mathbf{v} = \alpha_1 \mathbf{1}_{C_{n_1}} + \ldots + \alpha_k \mathbf{1}_{C_{n_k}}$ in
    $\mathbf{x}^\top (\alpha_1 \mathbf{1}_{C_{n_1}} + \ldots + \alpha_k \mathbf{1}_{C_{n_k}}) = 0$.
    Ker je to ena enačba z $k$ spremenljivkami, je dimenzija rešitvenega prostora enaka $k-1$.
    Posledično je $\dim\ker(A-2I) = k-1$.
\end{dokaz}


\begin{trditev}\label{tr:pet_ciklov}
    Podgraf $H[N(Q)]$ je enak disjunktni uniji petih ciklov $C_5$.
\end{trditev}
\begin{dokaz}
    Vzemimo graf $H$ z lastnimi vrednostmi 
    $\lambda_1 \geq \lambda_2 \geq \ldots \geq \lambda_{50}$ ter njegov podgraf 
    $H[v \cup N(Q)]$ z lastnimi vrednostmi 
    $\mu_1 \geq \mu_2 \geq \ldots \geq \mu_{26}$. Na njem uporabimo izrek o 
    prepletanju lastnih vrednosti~\ref{iz:interlacing}, s $r=26$ ter $n=50$.
    Za vsak $i \in \{2, 3, 4, 5\}$ dobimo:
    \[2 = \lambda_i \geq \mu_5 \geq \lambda_{24 + i} = 2.\]
    Večkratnost lastne vrednosti $2$ v $H[v \cup N(Q)]$ je torej vsaj $4$.

    Po lemi~\ref{le:knjiga}, uporabljeni na
    $G = H[N(Q)]$, $v \in Q$ ter $W = \gr_i$, dobimo, da je večkratnost lastne vrednosti $2$
    v $H[N(Q)]$ vsaj $4+1=5$, ker je $5$ tudi največje število ciklov, ki jih lahko 
    imamo v $20$-vozliščnem $2$-regularnem grafu, je 
    po lemi~\ref{le:cikellastne}, $H[N(Q)]$ 
    enak $C_5$.
\end{dokaz}






Velja torej $H[N(Q)] = 5C_5$. Označimo te cikle z $P_0, P_1, P_2, P_3, P_4$. Za vsak $i$ velja, da
$P_i \cup Q \subset H$ sestoji iz 10 vozlišč, je ožine 5, ter je kubičen graf (namreč 
vsako vozlišče iz $Q$ ima dva soseda v $Q$ ter natanko enega v $P_i$, simetrično velja za vozlišča iz $P_i$); 
to so karakteristike Petersenovega grafa, zato torej $P_i \cup Q$ inducira Petersenov graf. Ker bi našo analizo lahko začeli s poljubnim ciklom $C_5$,
bi lahko izbrali tudi $P_i$ namesto $Q$ ter tako izvedeli, da je $N(P_i) = Q \cup N_2(Q) = 5C_5$ 
(namreč vsako vozlišče iz $N_2(Q)$ ima natanko 5 sosedov v $N(Q)$, natanko enega na gručo). Posledično
je iz te perspektive razvidno, da je $H[N_2(Q)] = 4C_5$.
Označimo te cikle z $Q_1, Q_2, Q_3, Q_4$, originalni $Q$ pa s $Q_0$. Tudi $Q_i \cup P_j$ inducira Petersenov graf.
 
Postavimo $Q_i$ ter $P_j$ v matriko, ker poljubna unija tvori Petersenov graf, 
je jasno, lahko povezavo med $Q_i$ in $P_j$ opišemo z neko permutacijo 
$\pi_{ij}$. Oštevilčimo vozlišča v $P_j$ v smeri urinega kazalca, $Q_i$ pa 
najprej preoblikujemo v komplementarno zvezdno obliko, ter nato oštevilčimo prav 
tako v smeri urinega kazalca. BSŠ lahko predpostavimo $\pi_{0j} = \pi_{i0} = \text{id}$; identiteta
(namreč, sicer lahko cikle $Q_i$ ter $P_j$ preuredimo).
Opazimo, da za $i, j \leq 1$ permutacija $\pi_{ij}$ ne sme imeti fiksne točke, da se 
prepričamo vzemimo fiksno točko $k$ (torej $\pi_{ij}(k) = k$), potem nastane 
$C_4 = Q_0^k P_0^k P_j^k Q_i^k Q_0^k$, protislovje (z ožino). Dodatno, ker je $Q_i \cup P_j$ 
Petersenov graf, je $\pi_{ij}$ lahko zgolj rotacijska permutacija (namreč sicer nastanejo cikli), označimo 
jo z $\rho$, kjer je $\rho = (4\,0\,1\,2\,3)$ ter 
$\rho^m = (\rho^{m-1}(4)\, \rho^{m-1}(0)\, \rho^{m-1}(1)\, \rho^{m-1}(2)\, \rho^{m-1}(3))$,
kjer $\rho^0 = \id$.
Dodatno vidimo tudi, da si morajo biti permutacije, ki so v isti vrstici ali istem stolpcu,
različne. Za protislovje BSŠ vzemimo $\pi_{ij} = \pi_{i'j}$, tako nastane cikel
$Q_0^k Q_i^k P_j^{\pi_{ij}(k)} Q_i'^{k} Q_0^k$; protislovje.
Tako lahko poljubno izberemo različne si permutacije $\pi_{i1}$ in $\pi_{1j}$, dokler so si različne, BSŠ 
vzemimo kar naravno izbiro, $\pi_{i1} = \rho^i$ in $\pi_{1j} = \rho^j$. Opazimo še, da je sedaj lahko 
permutacija $\pi_{32}$ lahko bodisi $\rho$ ali $\rho^3$; vendar pri izbiri $\pi_{32} = \rho^4$ potem 
je $\pi_{11}^{-1} \pi_{12} \pi_{32}^{-1} \pi_{31} = \text{id}$, kar pomeni, da obstaja cikel $C_4$. Od tu 
naprej lahko brez težav zapolnimo preostale permutacije, saj so si vse permutacije v istem stolpcu ali vrstici 
različne.  

\begin{equation*}
\begin{array}{c|ccccc}
    & P_0 & P_1 & P_2 & P_3 & P_4 \\
    \hline 
    Q_0 & \text{id} & \text{id} & \text{id} & \text{id} & \text{id} \\
    Q_1 & \text{id} & \rho & \rho^2 & \rho^3 & \rho^4 \\
    Q_2 & \text{id} & \rho^2 & \rho^4 & \rho & \rho^3 \\
    Q_3 & \text{id} & \rho^3 & \rho & \rho^4 & \rho^2 \\
    Q_4 & \text{id} & \rho^4 & \rho^3 & \rho^2 & \rho \\
\end{array}
\end{equation*}

Tako smo konstruirali unikatnega kandidata za Mooreov graf stopnje $7$ in premera $2$,
ki ga imenujemo Hoffman-Singeltonov graf.
Preostane nam le še pokazati, da je graf resnično Mooreov in v zadnjih korakih nismo ustvarili
kratkih ciklov.

Videli smo, da $Q_i \cup P_j$ inducira Petersenov graf, ki ima ožino $5$. Torej, če bi 
cikel obstajal bi moral biti potekati vsaj skozi enega $P_i$-ja in dva $Q_i$-ja (ali pa obratno).
Ampak če cikel obišče enega $P_i$-ja in dva $Q_i$-ja, potem mora vsebovati vsaj dve vozlišči
iz vsakega, vendar potem je dolžina takega cikla vsaj $6$. Namreč $Q_i^u P_j^v$ lahko nadaljujemo 
edino z $Q_i^v$, kar pa niti ne tvori cikla, torej res mora vsebovati vsaj dve vozilšči iz vsakega.

Edina ostala možnost, ki bi lahko ustvarila cikel, je, da obišče dva $P_i$-ja in dva $Q_i$-ja,
in s tem tvori $C_4$. Recimo, da je ta cikel $Q_i^u P_j^v Q_k^w P_l^x Q_i^u$.

Vidimo lahko, da se permutacijo, ki poveže $Q_i$ z $P_j$ lahko zapišemo tudi krajše
kot $\pi_{ij} = \rho^{ij \pmod 5}$. Torej iz vozlišča $Q_i^u$ v $P_j^v$ lahko pridemo
le, če je $v = u + ij \pmod 5$, podobno se nazaj vračamo z inverzno permutacijo in zato
lahko iz vozlišča $P_j^v$ pridemo v $Q_k^w$ le, če je $w = v - jk \pmod 5$.

Tako bi naš cikel $Q_i^u P_j^v Q_k^w P_l^x Q_i^u$ lahko obstajal le, če bi veljale enačbe:
\begin{align}
    v &\equiv u + ij \pmod 5 \\
    w &\equiv v - jk \pmod 5 \\
    x &\equiv w + kl \pmod 5 \\
    u &\equiv x - li \pmod 5
\end{align}
Seštete enačbe dajo $v + w + x + u \equiv v + w + x + u + ij - jk + kl - li \pmod 5$, 
kar se poenostavi v $0 \equiv (i-k)(j-l) \pmod 5$. Ampak potem je bodisi $(i-k) \equiv 0$
bodisi $(j-l) \equiv 0$, ker so $i, j, k, l$ elementi množice $\{0, 1, 2, 3, 4\}$,  
to prisili, 
da je $i=k$ ali $j=l$, kar pomeni, da cikel ne obstaja.

Tako smo pokazali, da je graf, ki smo ga konstruirali resnično Mooreov graf stopnje $7$ in 
premera $2$, ter da je unikaten.

\begin{trditev}\label{tr:sosedstva_ciklov}
    Sosedstva med cikli $P_i$ in $Q_j$ so opisana z relacijo
    \begin{equation}
        P^i_j \sim Q^k_m \iff m \equiv j + i\cdot k \pmod 5
    \end{equation}
\end{trditev}

\paragraph{Grupa avtomorfizmov}
\begin{trditev}\label{tr:red_avtomorfizmov}
    Red avtomorfizmov Hoffman-Singeltonovega grafa je $|\aut(H)| = 50\cdot 7!$.
\end{trditev}

Fiksirali bomo nek cikel $C_5$ v grafu $H$, recimo $P^0$ in opazovali njegovo orbito
ter stabilizatorsko grupo. Če označimo $G = \aut(H)$, nam izrek o 
orbitah in stabilizatorjih~\ref{iz:orbita_stabilizator_red}
pove, da je $|\aut(H)| = |G\cdot P^0||G_{P^0}|$.

Število elementov v orbiti $G\cdot P^0$ je enako številu ciklov $C_5$ v grafu $H$,
v katere lahko preslikamo $P^0$.
Red stabilizatorja $G_{P^0}$ pa bomo izračunali tako, da bomo našli grupo $\Delta$, ki
ji je v njem vsebovana in pokazali, da mu je v resnici enaka.

\begin{lema}
    Število ciklov $C_5$ v grafu $H$ je enako $1260 = 50\cdot 7\cdot 6\cdot 6$
    In posledično je $|G\cdot P^0| = 50\cdot 7\cdot 6\cdot 6$.
\end{lema}
\begin{dokaz}
    Ker ima graf $H$ premer $2$ in ožino $5$, vsaka pot dolžine $4$ enolično določa
    cikel $C_5$. Tako pot lahko začnemo v poljubnem vozlišču, nadaljujemo pa
    jo z poljubnim sosedom, nato pa z vozliščem, ki ni sosed prvega, 
    ampak je sosed drugega, in tako naprej.
\end{dokaz}

\begin{lema}\label{le:permutira_qi_je}
    Vsak avtomorfizem, ki fiksira $P^0$ permutira $Q^i$-je, brez rotacij.
\end{lema}
\begin{dokaz}
    Naj bo $\sigma \in G_{P^0}$ in $Q = \{Q^0, \ldots, Q^4\}$. 
    Potem je $P^0_j \sim Q^k_j$ (\ref{tr:sosedstva_ciklov}).
    Tako za presekano množico sosedov velja $N(P^0_j)\cap Q = \{Q^k_j \mid k \in \mathbb{Z}_5\}$.
    Ker $\sigma$ ohranja sosedstva in fiksira $P^0_j$ velja 
    $\sigma(\{Q^k_j \mid k \in \mathbb{Z}_5\}) = \{Q^k_j \mid k \in \mathbb{Z}_5\}$
    (\ref{tr:stabilizator_sosedstvo}), kar pomeni, da $\sigma$ permutira $Q^i$-je brez rotacij.
\end{dokaz}

Definirajmo zdaj grupo $\Delta = \langle\alpha\beta\rangle$ generirano
z $\alpha$ in $\beta$, kjer sta $\alpha, \beta: V(H) \to V(H)$, definirani kot:
    \begin{align}
        &\alpha: \begin{cases}
            P_j^i &\mapsto P_j^{3i} \\
            Q_j^k &\mapsto P_j^{2k}
        \end{cases}\\
        &\beta: \begin{cases}
            P_j^i &\mapsto P_{j-i}^i \\
            Q_j^k &\mapsto Q_j^{k+1}
        \end{cases}
    \end{align}
\begin{lema}
    Grupa $\Delta$ je izomorna semidirektnemu produktu 
    (\ref{def:semidirektni_produkt}) ciklične grupe $C_5$ in grupe $C_4$.
    Posledično je $|\Delta| = 20$.
\end{lema}
\begin{dokaz}
    Najprej preverimo, da sta $\alpha$ in $\beta$ resnično avtomorfizma grafa $H$. 
    Po trditvi~\ref{tr:sosedstva_ciklov} je $P^i_j \sim Q^k_{j + ik}$.
    Res preslikavi ohranjata sosedstva, saj 
    \begin{align*}
    &\alpha(P^i_j) = P^{3i}_j \sim Q^{2i}_{j + (3i)(2k)} 
        \equiv Q^{2i}_{j + ik} = \alpha(Q^k_{j + ik}) \\
    &\beta(P^i_j) = P^i_{j-i} \sim Q^{k+1}_{(j-i)+i(k+1)} 
        \equiv Q^{k+1}_{j + ik} = \beta(Q^k_{j + ik})
    \end{align*}

    Dalje pokažimo, da je $\beta$ normalna podgrupa (\ref{def:normalna_podgrupa}) v $\Delta$.
    Velja namreč
    \begin{align*}
    &(\alpha\circ\beta\circ\alpha^{-1})(V(H)) \ni (\alpha\circ\beta\circ\alpha^{-1})(P^i_j) = (\alpha\circ\beta)(P^{2i}_j) = 
    \alpha(P^{2i}_{j-2i}) = P^{3(2i)}_{j-2i} \\
    &\equiv P^i_{j-2i} = \beta^2(P^i_j) \in \langle\beta\rangle \\
    &(\alpha\circ\beta\circ\alpha^{-1})(V(H)) \ni (\alpha\circ\beta\circ\alpha^{-1})(Q^k_j) = (\alpha\circ\beta)(P^{3k}_j) = 
    \alpha(P^{3k}_{j-3k}) = P^{3(3k)}_{j-3k} \\
    &\equiv P^k_{j-3k} = \beta^3(Q^k_j) \in \langle\beta\rangle
    \end{align*}

    Ker je tudi $\langle\beta\rangle \cap \langle\alpha\rangle = \{\id\}$,
    namreč $\alpha$ fiksira $Q^0$ medtem ko ga $\beta$ nikoli ne,
    sedaj vemo, da je $\Delta = \langle\beta\rangle \rtimes \langle\alpha\rangle$, 
    semidirektni produkt (\ref{def:semidirektni_produkt}), in da
    velja $|\Delta| = |\langle\beta\rangle||\langle\alpha\rangle|$ 
    (\ref{op:lastnost_semidirektnega_produkta}).

    Ker je $\beta^5 = \id$ in $\alpha^4 = \id$, je $|\langle\beta\rangle| = 5$ in 
    $|\langle\alpha\rangle| = 4$, posledično je $|\Delta| = 20$.

    Mimogrede, ker je $\alpha$ edini generator grupe $\langle\alpha\rangle$ in je reda
    $4$, je $\langle\alpha\rangle$ izomorfna ciklični grupi $C_4$. 
    Podobno je $\langle\beta\rangle$ izomorna ciklični grupi $C_5$. Tako tudi 
    $\Delta \cong C_5 \rtimes C_4$.
\end{dokaz}


\begin{lema}
    Grupa $\Delta$ je natanko stabilizatorska grupa $G_{P^0}$.
    Posledično je $|G_{P^0}| = 20$.
\end{lema}
\begin{dokaz}
    Če bi $G_{P^0}$ imel še kako permutacijo $\tau \not\in \Delta$,
    potem bi ta prisilila obstoj netrivialne permutacije, ki bi fiksirala $P^0$
    in še dva $Q^i$-ja. Namreč po Lagrangevem izreku~\ref{iz:index_lagrange} 
    vemo, da je $|G_{P^0}|$ večkratnik reda svoje podgrupe $\Delta$, torej večkratnik
    števila $20$. Če je grupa $G_{P^0}$ večja od grupe $\Delta$, 
    mora veljati $|G_{P^0}| \geq 40$. Po izreku o orbitah in 
    stabilizatorjih~\ref{iz:orbita_stabilizator_red} 
    je $|G_{P^0}| = |G_{P^0}\cdot Q^0||(G_{P^0})_{Q^0}|$.
    Ker $G_{P^0}$ fiksira $P^0$, permutira zgolj $Q^i$-je brez 
    rotacij (\ref{le:permutira_qi_je}), se $Q^0$ pri fiksiranem
    $P^0$ lahko preslika natanko v enega izmed $5$-ih $Q^i$-jev, 
    zato $|G_{P^0}\cdot Q^0| = 5$.

    Velja torej $|(G_{P^0})_{Q^0}| \geq 40 / 5 = 8$. Zdaj izrek o orbitah 
    in stabilizatorjih~\ref{iz:orbita_stabilizator_red} uporabimo še enkrat, 
    in tokrat dobimo 
    $|(G_{P^0})_{Q^0}| = |(G_{P^0})_{Q^0}\cdot Q_1||((G_{P^0})_{Q^0})_{Q_1}|$
    in posledično $|((G_{P^0})_{Q^0})_{Q^1}| \geq 8 / 4 = 2$.
    To nam neposredno pove, da obstaja vsaj ena netrivialna permutacija 
    $\sigma \in G_{P^0}$, ki fiksira $P^0$, $Q^0$ in $Q^1$. Označimo 
    $S = \{P^0, Q^0, Q^1\}$ in $\sigma(S) = S$. 

    Preostane nam še pokazati, da bi taka permutacija $\sigma$ nujno ustvarila
    cikel $C_4$, kar bi bilo protislovje.

    Vzemimo poljubno vozlišče $P^i_j$, kjer $i\neq 0$. Potem sta njegova soseda v množici 
    $S$ enaka $\{Q^0_j, Q^1_{j+1}\} = N(P^i_j)\cap S$ (\ref{tr:sosedstva_ciklov}). 
    Ker $\sigma$ fiksira $S$ 
    točkovno, velja $N(\sigma(P^i_j)) \cap S = N(P^i_j) \cap S$, %% TODO: Utemelji. z referenco.
    torej si $\sigma(P^i_j)$ in $P^i_j$ delita soseda $Q^0_j$ in $Q^1_{j+1}$.
    Tako smo prispeli v protislovje, našli smo namreč cikel
    $C_4 = P^i_j Q^0_j \sigma(P^i_j) Q^1_{j+1} P^i_j$. 
\end{dokaz}

Zdaj končno vidimo, da je $|\aut(H)| = |G\cdot P^0||G_{P^0}| = (50\cdot 7\cdot 6\cdot 6) \cdot 20 = 50 \cdot 7!$.
Naša trditev~\ref{tr:red_avtomorfizmov} je dokazana.

\begin{opomba}\label{op:hs_tranzitiven}
Vidimo lahko, da je $50\cdot 7!$ tudi največje število avtomorfizmov,
ki jih lahko ima graf reda $50$ in stopnje $7$. Namreč, po
izreku o orbitah in 
stabilizatorjih~\ref{iz:orbita_stabilizator_red}, 
je $|\aut(H)| = |G\cdot v||G_v|$ za poljubno vozlišče $v$.

Orbita vozlišča $v$ je množica vseh vozlišč, v katera ga lahko preslikamo in 
takih je največ $50$. Stabilizator $G_v$ pa lahko kvečjemu permutira vseh 
njegovih $7$ sosedov, kar pomeni, da je njegov red kvečjemu $7!$.

Ker je orbita poljubnega vozlišča $v$ maksimalna, je 
Hoffman-Singeltonov graf tudi vozliščno tranzitiven.
\end{opomba}

\begin{opomba}
    Izkaže se, da grupa avtomorfizmov Hoffman-Singeltonovega grafa
    izomorfna semidirektnemu produktu med projektivno posebno unitarno
    grupo $\text{PSU}(3, 5^2)$ in ciklično grupo $C_2$.
    Velja torej $\aut(H) \cong \text{PSU}(3, 5^2) \rtimes C_2$.
\end{opomba}





\subsubsection{Graf valence 57}
Obstoj Mooreovega grafa stopnje $\Delta = 57$ in premera $D = 2$ je še 
vedno odprto vprašanje. Vemo pa, da po definiciji Mooreovega grafa~\ref{def:moore},
tak graf, 
če obstaja $3250$ vozlišč. 
Po izreku Hoffman-Singeltona~\ref{iz:hoffmansingeltonizrek} pa tudi, 
da je stopnje $57$, da ima ožino $5$, ter
da so njegove lastne vrednosti $57$, $7$ in $-8$ z 
večkratnostmi $1$ ter $1729$ in $1520$.

Iz meje za neodvisno število $\alpha$ (moč neodvisne množice), ki jo 
podaja~\citeauthor{haemers2021hoffman}, ki pravi, da za $k$-regularni graf reda $n$ velja 
$\alpha \leq n \frac{-\lambda_{\min}}{\lambda_{\max} - \lambda_{\min}}$,
kjer je $\lambda_{\max}$ največja lastna vrednost, ter $\lambda_{\min}$ najmanjša lastna vrednost,
dobimo, da je neodvisno število tega grafa manjše ali enako $400$.

Predstavimo še trditev, ki pravi, 
da graf valence $57$, če obstaja, ne more biti vozliščno tranzitiven, 
kar pomeni, da ne more imeti simetrije, ki bi omogočala, da se katerokoli vozlišče preslika v 
katerokoli drugo vozlišče. To ga loči od ostalih treh Mooreovih grafov, ki so vsi vozlično tranzitivni.

\begin{trditev}
Graf valence $57$, če obstaja, ni vozlično tranzitiven.
Glej vira \citeauthor{higman1964finite} in \citeauthor{aschbacher1971nonexistence}. NEDOKONČANO.

\end{trditev}



\begin{dokaz}[Skica dokaza]
Predpostavimo, da je graf valence $57$ vozliščno tranzitiven. Potem njegova grupa avtomorfizmov
deluje tranzitivno na njegovih $3250$ vozliščih. Ker je premer grafa $2$, 
stabilizator $G_v$ razdeli preostalih $(n-1)$ vozlišč na natanko dve množici (orbiti).
Prva orbita vsebuje $57$ vozlišč, ki so sosedi vozlišča $v$,
druga orbita pa vsebuje $57\cdot (57 - 1) = 3192$ vozlišč, ki niso sosedna vozlišču $v$.
Skupaj s trivialno orbito $\{v\}$ to pomeni, da ima grupa rang $3$.

Predpostavimo sedaj, da je $G^\Omega$ grupa ranga $3$ stopnje $3250$ in podstopnje $57$.
Naj bo $\infty$ poljubna točka $\Omega$ in naj bo $H = G_{\infty}$ stabilizator te točke.
$H$ tako vsebuje vse avtomorfizme, ki fiksirajo točko $\infty$.

Naj bosta $\Delta$ in $\Gamma$ netrivialni orbiti grupe $H$ na množici $\Omega$ z $|\Delta| = 57$.
$\Delta$ je tako množica sosedov točke $\infty$, $\Gamma$ pa množica nesosedov točke $\infty$.

Naj bo zdaj $K$ podgrupa $H$ sestavljena iz vseh permutacij, ki so sode tako na $\Delta$ 
kot na $\Gamma$. Permutacije so sode, če zamenjajo sodo število elementov.
Rezultati iz (TODO?Section1) pokažejo, da involucija $\sigma$ na $H$ fiksira 
$1, 3, 7, 55$ ali $57$ elementov iz $\Delta$. Ker je $K$ soda podgrupa, so 
v $K$ možne le involucije, ki fiksirajo $1$ ali $57$ elementov iz $\Delta$. Namreč
zgolj $(57 - 1)/2 = 28$ in $(57 - 57)/2 = 0$ sta sodi števili.

Ker pa involucija $\sigma$ v $K$ poleg točk iz $\Delta$ fiksira tudi točko $\infty$, fiksira 
zvezdo s $55$???? ali $58$ točkami. Ker mora biti $\sigma$ soda tudi na $\Omega$, fiksira 
zvezdo s $58$ točkami.

TODO: Nadaljuj.
\end{dokaz}

\begin{opomba}
    \citeauthor{higman1964finite}\cite[str. 150]{higman1964finite} je pokazal, da
    so $C_5$, Petersenov graf in Hoffman-Singeltonov graf vsi vozliščno tranzitivni, 
    tako da je našel tranzitivne podgrupe njihovih grup avtomorfizmov.
    In sicer so te: Dihedralna grupa $D_10$ za $C_5$, simetrična grupa $S_5$ za Petersenov graf
    in projektivna specialna unitarna grupa $U_3(5)$ za Hoffman-Singeltonov graf.
\end{opomba}

\subsection{Neobstoj Mooreovih grafov}

%TODO
Glej \cite{bannai1973finite}.


\section{Razširitev Mooreovih grafov}\label{sec:moorerazsiritev}
% povezovalec, definicija, opomba pri definiciji da se ujema s prejšnjo

Za Mooreov graf velja, da ne le doseže maksimalno število vozlišč za dano maksimalno
stopnjo $\Delta$ ter premer $D$, temveč tudi doseže minimalno število vozlišč za dano
minimalno stopnjo $\Delta$ ter ožino $g=2D + 1$. Spodnjo mejo za število vozlišč za graf
z dano minimalno stopnjo $\Delta$ in ožino $g$ lahko naravno razširimo tudi za sode ožine.

Začnimo z dvema povezanima vozliščema $u$ in $v$, ko začnemo potovati iz vozlišča $u$ v eno smer
ter z vozlišča $v$ v drugo, se ti poti ne smeta srečati pred $k-1$ koraki -- namreč to bi vodilo 
v protislovje z izbrano ožino $g=2k$, ker pa minimiziramo število povezav predpostavimo, da se 
poti stakneta na $k$-tem koraku. Tako smo začeli z ničelnim nivojem, kjer smo imeli dve vozlišči $u$ 
in $v$, v prvem nivoju smo tako imeli $2(k - 1)$ novih vozlišč, saj je minimalna stopnja enaka $k$,
ter zaradi minimiziranja števila vozlišč je graf $k$ regularen. V $i$-tem nivoju imamo $2(k-1)^i$ 
vozlišč. Tako smo izpeljali spodnjo mejo enako 
$2\sum_{i=0}^{k-1} (k - 1)^i = 2\frac{(\Delta-1)^k - 1}{k - 2}$.

\begin{trditev}\label{tr:transition}
Naj bo $G$ regularen graf stopnje $\Delta$, ter z dolžino 
najkrajšega cikla, oziroma ožino, $g$. Tedaj obstaja spodnja meja 
za red grafa $G$ ter je enaka:
\begin{equation} %V opisu diplomske naloge je bila napaka.%
    |V(G)| \geq \begin{cases}
    \frac{\Delta(\Delta - 1)^{(g-1)/2}}{\Delta - 2}, & g \text{ je liha} \\
    \frac{2(\Delta - 1)^{g/2} - 2}{\Delta - 2}, & g \text{ je soda}
    \end{cases}
\end{equation}
\end{trditev}

V nadaljevanju vzemimo trditev~\ref{tr:transition} kot razširitev Moorevih grafov, pravimo 
da so tudi grafi, ki za dano ožino $g$ ter minimalno stopnjo $\Delta$ dosežejo minimalno število 
vozlišč, Mooreovi oziroma razširjeno Mooreovi. 

%1.
Mooreovi grafi sode ožine so dvodelni po konstrukciji, namreč lahko si predstavljamo, da smo pot iz 
vozlišča $u$ polagali izmenično (glede na nivoje (recimo ali so sodi ali lihi)) 
neko množico $A$ in $B$, pot iz vozlišča $v$ pa ravno obratno. Vozlišča iz zadnjega 
nivoja ne smejo biti povezana z vozlišči zadnjega nivoja iste množice, saj bi sicer ustvarili manjše 
cikle. Po konstrukciji so tudi regularni stopnje $d$.
% NUJNO: Vstavi sliko, ki to vizualizira

%2. !!!!!!!!!!!!
da poljubna dva elementa ene particije (npr. premici) nista
 hkrati incidentna z več kot enim elementom druge particije (točko).

%3. 
Njegov premer je enak $d$.

Te lastnosti pa se ujemajo z definicijo posplošenega večkotnika, ki jo poda 
\citeauthor{fresan2021packing}; ter tudi z definicijo šibkega (vključojoč 
degenerirane primere) posplošenega večkotnika
glede na terminologijo avtorja \citeauthor{van2012generalized}.
% Lahko bi začel: "se ujema z definicijo šibkega posplošenega večkotnika avtorja
% ... kar pa bomo za naše namene imenovali kar posplošeni večkotnik (zanimajo nast tudi
% degenerirani primeri)
\begin{definicija}
Posplošeni večkotnik je incidenčna struktura drugega ranga $(\PP, \LL, \II)$,
za katero velja:
\begin{enumerate}
    \item Obstaja naravno število $q\geq 1$ 
\end{enumerate}
\end{definicija}


\subsection{Posplošeni večkotniki}
% TODO: Nujno vključi 1.5.4 Lemma. iz Gheld knjige, ki govori o |L|, |P|, ...
% TODO: 1.7.1 Theorem (Feit & Higman [1964]). pokaže n=3 -> projektivne ravnine in n=12 -> (1,1)
% TODO: 1.7.2 Theorem. meje za s in t. (s <= t**2, ...)
% TODO: 1.7.7 Remark, govori o tem kako je del kvadrat pri projektivnih ravninah ...

% Naj bo L množica premic in P množica točk, posplošeni večkotnik je definiran kot ...  
% širše gledano gamma = (P, L, I) predstavlja geometrijo drugega ranga, ... 

% intuativna def: premer grafa je natanko n, ožina je natanko 2n, 
\begin{definicija}
    Preslikava $f: G \to G$ je endomorfizem, če ohranja sosednost.
    Torej, če velja:
    \[uv \in E(G) \implies f(u)f(v) \in E(f(G)).\]
\end{definicija}

\begin{definicija}
    Če točka $p \in \PP$ leži na premici $l \in \LL$ potem paru $\{p, l\}$ pravimo komora.
    Dve komori sta si sosednji, če se razlikujeta zgolj v enem elementu. Galerija je zaporedje
    prehodov med komorami. Galerija je minimalna če je najkrajša pot.  
\end{definicija}


\begin{izrek}[Feit-Higmanov izrek]
    Naj bo $S$ končni posplošeni $n$-kotnik z $s + 1$ točkami na vsaki premici 
    ter $t + 1$ premicami na vsaki točki. Potem je $s = t = 1$ in $S$ je običajni večkotnik, 
    ali pa je $n \in \{2,3,4,6,8,12\}$. Če sta $s, t > 1$ potem je $n \in \{2, 3, 4, 6, 8\}$,
    in če je $n\in\{3,6\}$ je $st$ kvadrat, in če $n=8$ je $2st$ kvadratno število.
    % običajni je mišljen kot navaden, sicer pa thick.
\end{izrek}
\begin{opomba}
    Izrek sta leta 1964 dokazala \textcite{feit1964nonexistence}, alternativni dokaz sta 
    kasneje podala tudi \textcite{kilmoyer1973theorem}. Velja tudi, da so posplošeni $3$~-kotniki
    natanko projektivne ravnine. 
    Iz izreka sledi tudi to, da posplošeni osem-kotniki niso regularni, saj 
    $2st$ ni kvadratno število, če sta $s$ in $t$ večja od $1$. 
    Sicer ($\sqrt{2ss} = \sqrt{2}s \notin \ZZ$).
\end{opomba}

% Mogoče še kaj o posplošenem digonu
\subsubsection{Projektivne ravnine}

\begin{definicija}[Projektivna ravnina]
    Projektivna ravnina je incidenčna struktura drugega ranga $(\PP, \LL, \II)$ 
    sestavljena iz množice točk $\PP$, množice premic $\LL$ ter incidenčne relacije
    $\II \subset \PP \cup \LL$, ki zadostuje naslednjim lastnostim:
    \begin{itemize}
        \item Za dve različni točki obstaja natanko ena premica, ki je incidenčna z obema.
        \item Za dve različni premici obstaja natnako ena točka, ki je z njima incidenčna.
        \item Obstajajo štiri take točke, da nobena premica ni incidenčna z več kot dvema izmed njih.
    \end{itemize}
\end{definicija}

\begin{izrek}[Bruck-Ryserjev izrek]
    Če je projektivna ravnina končna, obstaja število $N$, da vsaka premica vsebuje $N+1$ točk, 
    in je vsaka točka vsebovana v $N+1$ premicah. Če je $N \equiv 1$ ali $2 \mod 4$, ter če
    brezkvadratni del števila $N$ vsebuje vsaj en prafaktor oblike $4k + 3$, potem ne obstaja
    končna projektivna ravnina z $N+1$ točkami na premici \cite{bruck1949nonexistence}. Po 
    Fermatovem izreku o vsoti dveh kvadratov \cite{enwiki1319561489} je pogoj s prafaktorjem ekvivalenten temu,
    da je brezkvadratni del vsota dveh kvadratov.
\end{izrek}

\begin{definicija}[Galoisovo polje]
    Galoisovo polje je množica s končnim številom elementov $s$, v kateri veljajo operacije
    seštevanja, odštevanja, množenja in deljenja (razen z ničelnim elementom). 
    Označeno ga z $\GF(s)$. Obstaja natanko takrat ko je $s=p^n$ za poljubno praštevilo $p$ in 
    naravno število $n$. 
\end{definicija}

\begin{opomba}
    Če je $p$ praštevilo je $\GF(p) = \mathbb{Z}_p := \{0, 1, 2, \dots, p - 1\}$. 
    Dodatno je $\GF(p^n) = \{a_{n-1}x^{n-1} + \dots + a_1x + a_0 \mid a_i \in \GF(p)\}$ 
    množica polinomov, kjer množenje definiramo tako, da računamo v modulu nerazcepnega polinoma
    $P(x)$ stopnje $n$ (torej polinom, ki se ga v $\GF(p^n)$ ne da razbiti na produkt drugih).
    Tako velja $\GF(p^n) \cong \GF(p)[x] \,/\, \langle P(x) \rangle$.
\end{opomba}

\begin{definicija}
    TODO: Ta definicija povzroči zmedo, bolje definirati PG(k, K) bolj abstraktno in PG(2, K) bolj detajlno.
Projektivni prostor dimenzije $k$ ter reda $q=p^n$ označena s $\PG(k, q)$ obstaja za vse 
$k>2$ ter vse $n$ in vsa praštevila $p$. Namreč lahko jih generiramo s pomočjo Galloisovih polj,
kjer so točke $\PP = \{\langle \mathbf{x} \rangle \mid \mathbf{x} \in V \setminus \{0\}\}$, 
kjer je $V = \GF(q)^{k+1}$ vektorsko polje nad $\GF(q)$, ter 
$\langle \mathbf{x} \rangle = \{\lambda \mathbf{x} \mid \lambda \in \GF(q)\}$ (večkratnike 
vektorja $\mathbf{x}$ smatramo kot enake). Premice pa so množica hiperravnin 
$\LL = \{\mathbf{H}_{\mathbf{u} }\mid \mathbf{u} \in V \setminus \{0\}\}$, kjer je
$\mathbf{H}_{\mathbf{u}} = \{\langle \mathbf{x} \rangle \in \PP \mid \mathbf{u}\cdot\mathbf{x} = 0\}$.
Incidenčna relacija pa je nato definirana tako, da je točka $\langle \mathbf{x} \rangle$ 
incidenčna z premico (hiperravnino) $\mathbf{H}_{\mathbf{u}}$, če je njun skalarni produkt enak nič:
\[\II = \{(\langle \mathbf{x} \rangle, \mathbf{H}_{\mathbf{u}}) \in \PP \times \LL \mid
\sum_{i=0}^{k} u_i x_i = 0\}.\]
\end{definicija}

\begin{definicija}
    Projektivna ravnina je projektivni prostor dimenzije $k=2$.
\end{definicija}

\begin{trditev}
    Vsaka projektivna ravnina je posplošeni trikotnik reda $(s, s)$ za nek $s \in \ZZ_+$.
    Posledično je incidenčni graf vsake projektivne ravnine posplošeni mooreov graf stopnje
    $s + 1$ in ožine $6$.
\end{trditev}
% Lahko še kaj o ~Dekastrovih "grših" projektivnih ravninah.

\begin{definicija}[!!!!!!!]
    Projektivni prostor dimenzije $k$ nad poljem $\mathbb{K}$ je definiran kot:
    \[\PG(k, \mathbb{K}) = (V \setminus \{0\})/\sim,\]
    kjer je $V = \mathbb{K}^{k+1}$ vektorski prostor, relacija $\sim$ pa pravi, da
    sta vektorja $x$ in $y$ ekvivalentna, če obstaja $\lambda \in \mathbb{K} \setminus \{0\}$,
    da je $y = \lambda x$. Točke so $\PP = \{\langle x \rangle \mid x \in V \setminus \{0\}\}$, kjer 
    je $\langle x \rangle$ definiran kot prej. Premice pa so ravnine 
    $\LL = \{\langle \alpha x + \beta y \rangle \mid \alpha, \beta \in \mathbb{K} \textit{ neničelna.}\}$
    Incidenčna relacija je definirana glede na vsebovanost.
\end{definicija}

\begin{trditev}
    Vse končne projektivne ravnine temeljijo na $K = \mathbb{F}_{p^n}$.
    TODO: Glej Joseph Wedderburn
    TODO: Check!
\end{trditev}
% VSAKA projektivna ravnina ki je zgrajena iz vektorskega prostora nujno zadošča
% Desarguesovemu izreku.

\subsubsection{Posplošeni štirikotniki}
\begin{definicija}
    Skladno z definicijo avtorja~\citeauthor{van2012generalized} je šibki posplošeni 
    štirikotnik definiran kot geometrija $\Gamma = (\PP, \LL, \II)$, za katero velja:
    \begin{enumerate}
        \item Naj bo $\{p, L\}$ antizastava. Potem obstaja natanko ena zastava $\{q, M\}$,
        za katero velja $p \II M \II L \II q$.
        \item Vsaka točka leži na vsaj dveh, a ne vseh, premicah, in dualno, vsaka premica
        nosi vsaj dve, a ne vse, točke.
    \end{enumerate}
    Geometrija $\Gamma$ je posplošeni štirikotnik, 
    če drži prvi aksiom, je "thick" in vsebuje vsaj eno antizastavo.
\end{definicija}

% except for two particular families of generalized quadrangles, which however
% have nothing to do with symplectic spaces. Glej T_2(O) in T_3(O) konstrukcija iz ovalov.
% Ahrens-Szekeres AS(q) pa ni relevantna s not= t
% Tip GQ,"Red (s,t)",Obstoj,Opomba za graf
% W(q),"(q,q)",Vsi q=pn,Vedno obstaja.
% "Q(4,q)","(q,q)",Vsi q=pn,Dual W(q). Isti graf le za q=2n.
% T2​(O),"(q,q)",q=2n,Novi grafi za neklasične ovale.
% Kantor/Roman,"(q,q)",q=pn,Eksotični neklasični grafi.
% Wedderburnov izrek, ki pravi, da je vsak končni obseg nujno polje.
% g je (σ,1)-linearnost pomen: g(a*v, w) = a^sigma g(v, w) in g(v, w*b) = g(v, w)*b
% Wittov indeks w je dimenzija največjega podprostora, 
% ki je v celoti izotropen (to pomeni, da je vsak vektor v 
% tem podprostoru ortogonalen na vse ostale, vključno s samim seboj).
% w=2: Največji izotropni objekti so premice
!!!(bolj splošno): V projektivni ravnini se vsaki dve premici sekata. Da to preprečimo, uvedemo funkcijo $g(x, y)$,
ki bo določala kateri vektorji so si sorodni (zožala bo vektorski prostor) \cite{introtogeneralized}.

Naj bo $\mathbb{K}$ (komutativno) polje in obravnavajmo v $\PG(3, \mathbb{K})$, glede na 
neko izbrano bazo, simplektično polarnost $\tau$, ki preslika točko $(y_0, y_1, y_2, y_3)$
v ravnino z enabo
\[y1 X_0 - y_0 X_1 + y_3 X_2 - y_2 X_3 = 0.\]
Premico $L$ v $\PG{3, \mathbb{K}}$ imenujemo popolnoma izotropna, če velja $L^\tau = L$.
Definirajmo naslednjo geometrijo $W(\mathbb{K})$:
\begin{enumerate}
    \item Točke $W(\mathbb{K})$ so točke $\PG(3, \mathbb{K})$,
    \item Premice $W(\mathbb{K})$ so popolnoma izotropne glede na $\tau$.
\end{enumerate}
Pokazali bomo, da je $W(\mathbb{K})$ posplošeni štirikotnik, ter da se za $\mathbb{K} = \mathbb{F}_q$
ujema z $\GQ(q, q)$ (!!!malo preveri, nadaljuj) \cite[str. ~63]{van2012generalized}.

% 174, 240, 312



% str. 15, 52, 94, 63
\subsubsection{Posplošeni šestkotniki}

% 1.3.12 The smallest generalized hexagon

% posplošeni 8-kotniki pa niso regularni, in ne bodo mogli bit mooreovi!
% oziroma so lahko regularni (glej Ree-Titson graf), ampak niso s=t.
% prevedi Galoisovo polje v končno polje.


\subsection{Algebraična konstrukcija posplošenih večkotnikov}



\subsection{Obstoj razširjenih Mooreovih grafov sode ožine}

\section{Sorodni problemi}

\subsection{Mooreovi grafi z defektom}


\subsection{Kletke}




















%%%%%%%%%%%%%%%%%%%%%%%%%%%%%%%%%%%%
%%%%%%%%%%%%%%%%%%%%%%%%%%%%%%%%%%%%
%%%%%%%%%%%%%%%%%%%%%%%%%%%%%%%%%%%% OLD | 
%%%%%%%%%%%%%%%%%%%%%%%%%%%%%%%%%%%%
%%%%%%%%%%%%%%%%%%%%%%%%%%%%%%%%%%%%

\section{Uvod}

\begin{lema}\label{le:regularnilastne}
Naj bo $G$ $k$-regularen graf in $A$ njegova sosednostna matrika. Potem je $k$ njegova lastna
vrednost in vse ostale lastne vrednosti ležijo na intervalu $[-k, k]$ in so si različne. 
Velja celo to, da je
$k$ enojna lastna vrednost, vendar dokaz tega izpustimo.
\end{lema}
\begin{dokaz}
    Ker je $G$ $k$-regularen, je vsako njegovo vozlišče povezano z $k$ sosedi. Posledično
    se vsaka vrstica sosednostne matrike $A$ sešteje v $k$. 
    Velja torej $A\mathbf{1} = k\mathbf{1}$, $k$ je lastna vrednost.

    Naj bo zdaj $\lambda$ poljubna lastna vrednost matrike $A$, ter $\mathbf{x}$ njen 
    pripadajoč lastni vektor. Naj bo $i^{*}$ indeks, za katerega velja $|x_{i^{*}}| = \max_{j} |x_j|$.
    Potem velja:
    \begin{align*}
        |\lambda| \cdot |x_{i^{*}}| &= |\lambda x_{i^{*}}| = |(A\mathbf{x})_{i^{*}}| = \left|\sum_{j} A_{i^{*}j} x_j\right| \leq \sum_{j} A_{i^{*}j} |x_j|
        \leq \sum_{j} A_{i^{*}j} |x_{i^{*}}| = k |x_{i^{*}}|.
    \end{align*}
    Posledično je $|\lambda| \leq k$.
\end{dokaz}

\section{Uvodne definicije}

\begin{definicija}
Graf $G$ je krepko regularen s parametri $(n, k, \lambda, \mu)$, če je
$G$ regularen stopnje $k$ na $n$ vozliščih, ter če imata poljubni dve sosedni 
vozlišči natanko $\lambda$ skupnih sosedov, medtem ko imata poljubni dve 
nesosednji vozlišči natanko $\mu$ skupnih sosedov.
\end{definicija}

\begin{definicija}
    Naj bosta $(G, \cdot)$ in $(H, \circ)$ grupi. 
    Preslikava $h: G \to H$ je homomorfizem, če za vse $u$ in $v$ v $G$ velja
    $h(u \cdot v) = h(u) \circ h(v)$. Od tod sledi, 
    da $h(e_G) = e_H$, kjer je $e_G$ enota grupe $G$, ter $e_H$ enota grupe $H$ ter
    da $h(u^{-1}) = h(u)^{-1}$ za vsak $u \in G$.
    Če je $h$ bijektivna, potem pravimo, da je $h$ izomorfizem.
    Dve grupi $(G, \cdot)$ in $(H, \circ)$  sta izomorfni, če obstaja izomorfizem
    med njima, pišemo $(G, \cdot) \cong (H, \circ)$.
\end{definicija}
\begin{definicija}
Avtomorfizem grafa $G$ je bijektivna preslikava $\sigma: V(G) \to V(G)$, ki ohranja sosednost, 
torej $uv \in E(G) \implies \sigma(u)\sigma(v) \in E(G)$, oziroma v drugi notaciji 
$u^\sigma v^sigma \in E(G)$. Torej avtomorfizem je izomorfizem grafa $G$ v samega sebe.
Kompozitum dveh avtomorfizmov je zopet avtomorfizem, zato avtomorfizmi grafa $G$,
pod operacijo kompozituma, tvorijo grupo, ki ji pravimo grupa avtomorfizmov 
grafa $G$ in jo označujemo z $\aut(G)$.
V grafu z $n$ vozlišči je $\aut(G) \subset S_n$, kjer je $S_n$ množica vseh permutacij $n$ elementov,
tako velja tudi $|\aut(G)| \leq n!$.
Naj bo dodatno $s \in \aut(G)$ tak, da $s^k = \id$ za nek $k \in \ZZ_+$. 
Potem pravimo, da je $s$ avtomorfizem reda $k$. Posebno, če $k=2$ potem pravimo,
da je $s$ involucija.
\end{definicija}

\begin{definicija}
Graf $G$ je vozliščno tranzitiven, če je za poljubni vozlišči $u$ in $v$ obstaja avtomorfizem
$\sigma \in \aut(G)$, da $\sigma(u) = v$. Torej grupa avtomorfizmov deluje tranzitivno na množici
vozlišč grafa $G$.
\end{definicija}



\begin{definicija}
Naj bo $G = (V, E)$ graf, kjer je $V$ množica njegovih
vozlišč, ter $E$ množica njegovih povezav. Red grafa $G$ je 
število njegovih vozlišč, torej $\red(G) = |V|$.
\end{definicija}

\begin{definicija}
Naj bo $G = (V, E)$ graf, ter $v \in V$ njegovo vozlišče. 
Stopnja vozlišča $v$ je enaka številu njegovih sosedov, torej
$\deg(v) := \#\{u; \quad u \sim v \land u \in E\}$. Stopnjo 
grafa $G$ pa označimo z maksimalno stopnjo njegovih vozlišč:
$\Delta(G) := \max_{v \in V(G)} \deg(v)$
\end{definicija}

\begin{definicija}
Premer grafa $G$ je največja razdalja med katerima koli 
vozliščema v grafu. Torej, $\diam(G) = \max_{u, v \in V(G)} d(u, v)$, 
kjer je $d : V(G) \times V(G) \rightarrow \N$ metrika, ki meri  
razdaljo med danima vozliščema. 
\end{definicija}

\begin{definicija}
Graf je $\Delta$ regularen, takrat ko je vsako njegovo vozlišče 
stopnje $\Delta$, torej $\forall v \in V(G): \quad \deg(v) = \Delta$.
\end{definicija}

\begin{definicija}
Graf $G = (V, E)$, reda $n$ je krepko regularen natanko tedaj ko je $k$ regularen 
in obstajata 
konstanti $\lambda$ in $\mu$, da imata vsaki sosednji vozlišči 
natanko $\lambda$ skupnih sosedov, vsaki  ne sosedni pa natanko 
$\mu$ skupnih sosedov. Krepko regularen graf z danimi prametri 
označimo s $\srg(n, k, \lambda, \mu)$.
\end{definicija}

\begin{definicija}
Ožina grafa $G$ je enaka dolžini najkrajšega cikla v $G$. Če je graf 
$G$ brez ciklov, označimo ožino grafa z neskončnostjo.
\end{definicija}

\begin{definicija}
Naj bo $G$ povezan acikličen graf na $n$ vozliščih ter $n - 1$ povezavah,
ki tvorijo zaporedje $v_1 \sim v_2 \sim \ldots \sim v_n$. Tedaj graf tvori 
pot od $v_1$ do $v_n$, označimo ga z $P_n$.
\end{definicija}

\begin{definicija}
Naj bo graf $G$ cikel na $n$ vozliščih. Natančneje naj bo $G$ graf $P_n$
s dodano povezavo $v_1 \sim v_n$. Tedaj tak graf označimo z $C_n$.
\end{definicija}

\begin{definicija}
Naj bo $G$ graf reda $n$ z množico vozlišč $V$ ter množico povezav $E$. Tedaj je 
njegova matrika sosednosti $A \in \ZZ ^ {n \times n}$ zadostuje:
\begin{equation}
A_{u, v} := \left(\# \textit{povezav dolžine 1 med vozliščema $u$ in $v$}\right)_{u, v}.
\end{equation}
\end{definicija}

\begin{definicija}
Naj bosta $u$ in $v$ vozliči v danem grafu. Pravimo, da je vozlišče 
$v$ v dosegu vozlišča $u$, če v grafu obstaja pot od $u$ do $v$, takrat 
označimo $v \in \reach(u)$.
\end{definicija}

\begin{definicija}
Naj bo $G$ graf in $H \subset G$ njegov podgraf. $N_k(H)$ označuje 
množico vozlišč, kjer je vsaka od $H$ oddaljena za $k$ in v $H$ ni 
vsebovana.
\end{definicija}

\begin{izrek}\label{iz:polinom_v_matriki}
Naj bo $A$ kvadratna matrika dimenzije $n$, $v$ njen lastni 
vektor za lastno vrednost $\lambda$, $p$ polinomska funkcija,
potem velja $p(A)v = p(\lambda)v$.
\end{izrek}

\begin{dokaz}
Velja namreč:
\begin{equation}
p(A)v = \left(\sum_{k} c_k A^k\right)v = 
\sum_{k} c_k \lambda^k v = p(\lambda)v
\end{equation}
\end{dokaz}

\begin{izrek} [Poincaré separation theorem]
...
\end{izrek}

%%%%%%%%%%%%%%%%%%%%%%%%%%%%%%%%%%%%%%%%%%%%%%%%%%%%%%%%%%%%%%%

\section{Moorejev graf}\label{sec:moorejev_graf}

Mooreovi grafi so ekstremalni grafi, ki glede na dan premer ter dano 
maksimalno stopnjo dosežejo maksimalno število vozlišč. Leta 1960 sta jih tako 
poimenovala Hoffman in Singelton po Edwardu F. Mooru, ki je prvi zastavil vprašanje o 
opisu in klasifikaciji teh struktur.

\begin{definicija}
Naj bo $G = (V, E)$ graf s premerom $D$ ter maksimalno stopnjo $\Delta$.
Če je njegov red enak
\begin{equation}
    |V| = 1 + \Delta + \Delta (\Delta - 1) + \ldots + \Delta (\Delta - 1)^{D - 1},
\end{equation}
potem pravimo, da je $G$ Moorejev graf.
\end{definicija}

\begin{opomba}
Red Moorejevega grafa parametrov $(\Delta, D)$ je večji ali enak redu 
poljubnega grafa $P$, ki zadostuje maksimalni stopnji $\Delta$ ter premeru $D$.
Za poljubno vozlišče $v$ grafa $P$ namreč velja, da je število vozlišč od $v$ 
oddaljenih za $i$ največ toliko kot, če bi bilo vsako vozlišče razdalje $i-1$ 
od $v$ povezano z nadaljnimi $\Delta - 1$ vozlišči, oziroma
$\forall\; 0 \leq i \leq D:\;|N_{i}(v)| \leq |N_{i-1}(v)|\cdot (\Delta - 1)$. Induktivno
tako dobimo zgornjo mejo za red grafa $P$, ki je natanko enaka redu Moorejevega 
grafa.
\end{opomba}

\begin{trditev}\label{tr:regularnost_in_omejitve_dolzine}
Naj bo $G$ Moorejev graf s parametroma $(\Delta, D)$. Tedaj je
$G$ regularen stopnje $\Delta$, ter za poljubni dve vozlišči $u$ in $v$ za 
kateri $u \neq v$ velja, da med njima obstaja natanko ena najkrajša pot 
dolžine največ $D$.
\end{trditev}

\begin{dokaz}
Vzemimo poljubno vozlišče $u$ grafa $G$, $\deg(u)$ stopnja vozlišča $u$,
$\Delta$ pa maksimalna stopnja v grafu $G$. Potem analogno konstrukciji 
Moorejeve spodnje meje velja $|V(G)| \leq 1 + \deg(u)(\Delta - 1) + 
\ldots + \deg(u) (\Delta - 1)^{D - 1}$. Ker pa je Moorejeva meja v grafu
$G$ že dosežena, ter $\deg(u) \leq \Delta$, sledi, da je $\deg(u) = \Delta$.
Podobno se o enoličnosti poti dolžin $D$ prepičamo tako, da opazimo, da 
konstrukcija Moorejeve meje predpostavi, da ima vsako vozlišče $z$ v $i$-tem 
nivoju $N_i(z)$ natanko $\Delta - 1$ sosedov, ki niso povezani z nobenim drugim 
vozliščem iz $i$-tega nivoja, razen ko $i=D$. Ob zanemaritvi notranjih povezav 
zadnjega nivoja $N_D(z)$ skupek nivojev inducira drevesni podgraf za $G$, na katerem 
je vsaka pot med dvema vozliščema dolžine $D$ enolična.
\end{dokaz}

\begin{posledica}\label{po:ozina}
Po trditvi~\ref{tr:regularnost_in_omejitve_dolzine} v grafu s parametroma $(\Delta, D)$ 
med poljubnima vozliščema
$v$ ter $u$ obstaja največ ena najkrajša pot dolžine največ $D$. Posledično 
obstaja tudi pot od $u$ do $v$ dolžine $D + 1$. Najkrajši cikel 
v grafu je tako lahko napet na najmanj $D + (D + 1)$ vozlišč.
Ožina Moorejevega grafa je torej enaka $g=2D + 1$. !!!! utemeljitev ni vredu.

V nadaljevanju bomo zasledili, da so Mooreovi grafi deležni številnih omejitev, 
kar bo utemeljilo tudi njihovo redkost, namreč Mooreovi grafi obstajajo le za:
\begin{align}
\Delta = 1, D = 0; \text{graf $K_1$} \\
&\Delta = 2, D \geq 1; \text{cikli $C_{2D + 1}$} \\
&D = 2, \Delta \in \{1, 3, 7, 57\}; \text{grafi $P_2$, Petersenov graf, Hoffman-Singeltonov graf, neznani graf} \\
\end{align}

\end{posledica}

%%%%%%%%%%%%%%%%%

\subsection{Omejitve lastnih vrednosti}

Primer v katerem je premer $D = 1$ je trivialen, saj iz tega sledi,
da je tudi $\Delta = 1$ ter edini Moorejev graf s takimi parametri enak 
$P_2$. V nadaljnem torej predpostavimo $D, \Delta > 1$.

Tu bi omenil splošen problem in splošno enačbo ... 
\begin{trditev}
Mooreevi grafi zadostujejo naslednji matrični enačbi:
\begin{equation}
\begin{cases} 
G_{\Delta, 0} = I \\
G_{\Delta, 1} = A + I \\
G_{\Delta, i} = A G_{\Delta, i - 1} - (\Delta - 1) G_{\Delta, i - 2} & \text{za } i \geq 2
\end{cases}
\end{equation}
\end{trditev}

\begin{dokaz}
    Bazna primera za $i = 0$ in $i = 1$ sta očitna, saj $G_{\Delta, 0}$ meri število poti
    dolžine 0, torej je enaka identiteti, $G_{\Delta, 1}$ pa meri število poti dolžine
    največ 1, torej je enaka $A + I$. Za indukcijski korak pa opazujmo, da 
    $A G_{\Delta, i - 1}$ meri število poti dolžine največ $i$ med vozliščema $u$ in $v$,
    vključujoč tudi tiste, ki se iz $i-1$ nivoja vrnejo nazaj v $i-2$-gi nivo. 
    Ker pa je graf $G$ regularen stopnje 
    $\Delta$ za vsako vozlišče obstaja natanko $(\Delta - 1) G_{\Delta, i -2}$ 
    takih sprehodov, ki jih ustrezno odštejemo.
\end{dokaz}

\begin{opomba}
    V posebnem primeru za premer $D=2$, Mooreev graf zadođča enačbi:
    \[A^2 + A - (\Delta - 1) I = J,\]
    kjer je $J$ matrika samih enic.
\end{opomba}


\begin{trditev}
Za $A^k = (A^k_{u, v})_{u, v}$ velja:
\begin{equation}
A^k_{u, v} = \# \text{povezav dolžine $k$ med vozliščema $u$ in $v$}
\label{eq:potence_matrike_sosednosti}
\end{equation}
\end{trditev}

\begin{dokaz}
Vzemimo poljubni vozlišči $u$ in $v$ v grafu $G$. Ker velja 
\begin{align}
A^k_{u, v} &= \sum_{z \in V(G) \setminus \{u, v\}} A_{u, z} A^{k - 1}_{z, v} \\
           &= \sum_{z:~ z \sim u \land v \in \reach(u)} A_{u, z} A^{k - 1}_{z, v}
\end{align}
po definiciji matrike sosednosti $A$ ter indukcijski predpostavki res 
štejemo vse povezave dolžine $k$, ki preko vozlišča $z$ vodijo iz $u$ v $v$.
\end{dokaz}

Omejimo se sedaj na problem premera $D=2$ ter stopnje $\Delta > 1$. Predpostavimo,
da dosega Moorejevo mejo, torej je njegov red enak $\Delta^2 + 1$.
(Kar je tudi dimenzija njegove matrike sosednosti).
Ker je $G$ regularen stopnje $\Delta$ in ker med vsakima disjunktnima 
vozliščema obstaja ena sama pot dolžine največ $D$, za njegovo matriko sosednosti 
$A$ velja naslednja enačba: 
\begin{equation}
A^2 + A - (\Delta - 1) I = J
\end{equation}
, kjer je $I$ identiteta in $J$ matrika samih enic.
Matrika $A^2 + A$ namreč meri število poti dolžin največ $D$
\eqref{eq:potence_matrike_sosednosti}. Ker pa je graf 
regularen za vsako vozlišče obstaja natanko $\Delta$ poti dolžine $D$, ki se 
vrnejo nazaj, zato v enačbi odštejemo ustrezen faktor identične matrike.

Preko dane matrične enačbe lahko izračunamo spektralno polje incidenčne 
matrike $A$. Ker je karakteristični polinom za matriko $J$ enak 
$x^{\Delta^2}(x - (\Delta^2 + 1))$, sledijo naslednje enačbe:
\begin{align}
x^2 + x - (\Delta - 1) = \Delta^2 + 1 \label{eq:karakteristicna_1}\\ 
x^2 + 1 - (\Delta - 1) = 0 \label{eq:karakteristicna_2}
\end{align} 
Od koder dobimo lastne vrednosti $\Delta$ (po \eqref{eq:karakteristicna_1})
večkratnosti 1, ter lastni vrednosti $r$ in $s$ enaki 
$(-1 \pm \sqrt{4\Delta - 3})/2$. Več o večkratnosti lastnih vrednosti 
$r$ in $s$ pa nam podata enačbi $a + b = \Delta^2$ ter sled matrike, ki 
je enaka seštevku vseh lastnih vrednosti, po drugi strani pa 
je sled matrike sosednosti očitno enaka 0, torej: 
$\Delta + a r + b s = 0$. 

Če sta $r$ in $s$ iracionalni števili je $a = b$ ter $\Delta = 2$, 
sicer pa sta $r$ in $s$ celi števili, torej $4\Delta - 3 = m^2$ za nek 
$m \in \ZZ_+$. Ob substitucijah $a = \Delta^2 - b$ in 
$\Delta = \frac{m^2 + 3}{4}$ sledi:
\begin{align}
\Delta + a r + b s &= 0 \\
\Leftrightarrow \\
\frac{m^2 + 3}{4} + \left(\left(\frac{m^2 + 3}{4}\right)^2 - b\right) 
\frac{-1 + m}{2} + b \frac{-1 - m}{2} &= 0 \\
\Leftrightarrow \\
\frac{m^5}{32} + \ldots + \left(\frac{3}{4} - \frac{9}{32}\right)  &= 0\\
\Leftrightarrow \\
m^5 + \ldots + 15 &= 0
\end{align}

Preden nadaljujemo si oglejmo pomemben rezultat linearne algebre.

\begin{izrek}[Izrek o celoštevilskih ničlah]
Naj bo $p = a_n x^n + \ldots + a_0$ polinom z celimi koeficienti, kjer 
je vodilni koeficient neničelen. Če ima polinom celoštevilsko ničlo $q$,
potem $q$ deli $a_0$.
\end{izrek}

\begin{dokaz}
Ker je $q$ ničla polinoma $p$ velja $p(q) = a_n q^n + \ldots + 
a_0 = 0$. Po preureditvi dobimo $a_0 = -(a_n q^n + \ldots + a_1 q)$.
Ker je desna stran deljiva z $q$ je z $q$ deljiv tudi $a_0$.
\end{dokaz}

Ker ima enačba $x^2 + 1 - (\Delta - 1) = 0$ cele koeficiente so njene 
racionalne rešitve nujno cele (po izreku).
S tem rezultatom, ter z upoštevanjem, da je $m$ pozitivno celo število, 
lahko sedaj razberemo, da so možne rešitve za $m$ enake 
$\{1, 3, 5, 15\}$ kar so tudi deljitelji števila 15. Posledično 
so kandidati za $\Delta$ enaki $\{1, 3, 7, 57\}$. (Utemelji še, da
so koeficienti naračunane enačbe v $m$ res celi, in smo smeli 
uporabiti izrek.)

\begin{izrek}
Naj bo $G$ Moorejev graf, tedaj je njegova ožina enaka $g=$
\end{izrek}

%%%%%%%%%%%%%%%%%%%%%%%%%%%%%%%%%%%%%%%%%%%%%%%%%%%%%%%%%%%%%%%%%%%%

\section{Konstrukcija znanih Moorejevih grafov}

\subsection{Konstrukcija Petersenovega grafa}

\subsection{Poiskusi dokaza o obstoju oziroma neobstoju Moorejevega grafa valence 57}

\subsection{Dokazi o neobstoju ostalih Mooreevih grafov}

%%%%%%%%%%%%%%%%

\section{Posplošitev Moorejeve meje}
% bolje bi bilo: ... incidenčna relacija,da se reče kletka, ... kar pa je ravno definicija posplošenega večkotnika
% oglejmo si še enakovredno definicijo z geometrijo, omenimo znane rezultate, napeljemo na projektivno ravnino, ...
Podobno kot smo v sekciji~\ref{sec:moorejev_graf} izpeljali 
zgornjo mejo za red povezanih grafov s premerom $D$ ter 
maksimalno stopnjo $\Delta$, lahko izpeljemo tudi spodnjo mejo 
za soroden problem.

\begin{trditev}
Naj bo $G$ regularen graf stopnje $\Delta$, ter z dolžino 
najkrajšega cikla, oziroma ožino, $g$. Tedaj obstaja spodnja meja 
za red grafa $G$ ter je enaka:
\begin{equation} %V opisu diplomske naloge je bila napaka.%
    |V(G)| \geq \begin{cases}
    \frac{\Delta(\Delta - 1)^{(g-1)/2}}{\Delta - 2}, & g \equiv 1 \pmod 2\\
    \frac{2(\Delta - 1)^{g/2} - 2}{\Delta - 1}, & sicer
    \end{cases}
\end{equation}
\end{trditev}

\begin{dokaz}
    
\end{dokaz}

\begin{opomba}
Izkaže se, da so $\Delta$-regularni grafi
z ožino $g$, kateri izpolnijo spodnjo mejo za število vozlišč,
natanko Moorejevi grafi za lihe $g$. Za sode $g$ pa so ravno 
enaki incidenčnim grafom posplošenih večkotnikov reda $\Delta -1$,
o katerih bomo razglabljali v naslednjih poglavjih. Dodatno lahko pripomnimo, 
da grafom, ki dosežejo posplošeno Mooreevo mejo, pravimo \text{cage}, oziroma kletka, 
splošneje $(k g)$-kletka je graf, ki doseže minimalno število vozlišč, med vsemi grafi z 
maksimalno stopnjo $k$ in ožino $g$.
\end{opomba}

\begin{dokaz}
    
\end{dokaz}

%%%%%%%%%%%%%%%%%%%%%

\section{Geometrije}
\begin{definicija}
Geometrija (ranga n) je nabor
$\Gamma = \left(
    \A_1, \A_2, \ldots, \A_n, \II
\right)$, kjer so $\A_1, \A_2, \ldots, \A_n$ neprazne paroma disjunktne 
množice, in $\II$ relacija incidenčnosti za katero velja 
\[\II \subseteq \bigcup_{i\neq j} \A_i \times \A_j.\] Če za poljubna
$a, b \in \bigcup_{i} \A_i$ velja $(a, b) \in \II$, potem pišemo tudi
$\IIw{a}{b}$. 

Če za neko pozitivno celo število $m$ velja še 
\[\forall i \; \forall a \in \A_i\; \forall j \neq i : \ 
m = \#\{b \in \A_j\mid \IIw{a}{b} \},\]
potem pravimo, da je geometrija reda $m$.
\end{definicija}

\begin{definicija}
Naj bo $\Gamma = \left(
    \A_1, \A_2, \ldots, \A_n, \II
\right)$ geometrija. Pravimo, da je $\Gamma^D$ njena dualna geometrija 
če je $\Gamma^D = \left(
    \A_1, \A_2, \ldots, \A_n, \II^D
\right)$, kjer je $\IIwd{a}{b} \Leftrightarrow \IIw{b}{a}$.
\newline 
Podobno pravimo, da je $\Gamma'$ podgeometrija natanko 
tedaj ko $\Gamma' = \left(\A_1', \A_2', \ldots, \A_n', \II'\right)$,
kjer za vsak $1 \leq i \leq n$ velja $\A_i' \subseteq \A_i$ ter kjer 
je $\II'$ ustrezna zožitev relacije $\II$.
\end{definicija}



Navkljub splošnejšim definicijam si podrobneje poglejmo geometrije ranga 2
in strukture, ki le tam nastanejo.

\begin{definicija}
Geometrija (ranga 2) je nabor 
$\Gamma = \left(\PP, \LL, \II\right)$. Množici $\PP$ pravimo \textsl{točke}, 
množico $\LL$ pa poimenujmo \textsl{črte}. Pravimo, da je $p \in \PP$ 
\textsl{točka} ter $l \in \LL$ \textsl{črta}. \newline 
\textsl{Črta} je v relaciji z \textsl{točko}, če velja $\IIw{l}{p}$.
\end{definicija}

\begin{definicija}[Incidenčni graf]
Naj bo $\Gamma = \left(
    \A_1, \A_2, \ldots, \A_n, \II
\right)$ geometrija. Incidenčni graf geometrije $\Gamma$ je graf $G$
za katerega velja $V(G) = \bigcup_i \A_i$ ter 
$E(G) = \{uv \mid \IIw{u}{v} \lor \IIw{v}{u}\}$. Tak graf je očitno $n$-delen.
\end{definicija}

\subsection{Posplošeni večkotniki}

\begin{definicija}[Večkotnik]
Naj bo $\Gamma = (P, \mathcal{L}, I)$ geometrija ranga $2$. 
Pravimo, da v geometriji $\Gamma$ \emph{obstaja večkotnik dolžine $n$} 
natanko tedaj, ko obstajajo različni elementi 
\[
x_1, x_2, \ldots, x_{2n} \in P \cup \mathcal{L}
\]
za katere velja:
\begin{enumerate}
  \item elementi se izmenjujejo po vrstah, torej $x_{2i-1} \in P$ in $x_{2i} \in \mathcal{L}$ za vsak $i$;
  \item za vsak $i = 1, 2, \ldots, 2n$ velja $x_i \, I \, x_{i+1}$, pri čemer privzamemo $x_{2n+1} = x_1$.
\end{enumerate}
Zaporedju $(x_1, x_2, \ldots, x_{2n})$ pravimo \emph{večkotnik dolžine $n$} 
v geometriji $\Gamma$.
%
Opozarjamo, da obstoj večkotnika v geometriji ni nujno
ekvivalenten obstoju podgeometrije $\Gamma'$, katere incidenčni
graf je izomorfen ciklu $C_n$.
\end{definicija}

\begin{definicija}[Posplošeni večkotnik]
    Naj bo $n$ pozitivno celo število. Šibki posplošeni $n$-kotnik je geometrija
    $\Gamma = (\PP, \LL, I)$ ranga 2, ki zadošča naslednjima pogojema:
    \begin{enumerate}
        \item $\Gamma$ ne vsebuje običajnega $k$~-kotnika (kot podgeometrijo) za $0 \leq k < n$.
        \item Poljubna dva elementa iz $\PP \cup \LL$ sta vsebovana v nekem običajnem $n$~-kotniku
        (kot podgeometrija) v $\Gamma$, \emph{amartmaju}.
    \end{enumerate} 
    Posplošeni $n$~-kotnik je šibki posplošeni $n$~-kotnik $\Gamma$, ki zadošča še dodatnemu pogoju:
    \begin{enumerate}[resume]
        \item V $\Gamma$ obstaja običajen $(n+1)$~-kotnik (kot podgeometrija).
    \end{enumerate}, 
\end{definicija}

\begin{opomba}
    Vsak običajen $n$~-kotnik je tudi šibki posplošen $n$~-kotnik. % 1.3.10 Van Maldeghem.
\end{opomba}
Če posplošene večkotnike raje gledamo z vidika incidenčnih grafov, potem 
lahko zgornjo definicijo preoblikujemo na sledeč način.

\begin{definicija}\label{def:posploseniveckotnikiincidencnigraf}
    Naj bo $n$ pozitivno celo število. Incidenčni graf šibkih 
    posplošenih $n$~-kotnikov zadošča naslednjima pogojema: 
    \begin{enumerate}
        \item Ožina grafa je natanko $2n$.
        \item Premer grafa je natanko $n$.
    \end{enumerate}
    Graf $\Gamma$ je incidenčni graf posplošenega $n$~-kotnika, če zadošča še dodatnemu pogoju:
    \begin{enumerate}[resume]
        \item Stopnja vsakega vozlišča v $\Gamma$ je vsaj 3.
    \end{enumerate}
\end{definicija}

% lahko bi tu omenil še Theorem 1.3.8 Let (V, X), V = VI U V2, be a spherical building of rank 2. 
% Then r = (VI, V 2 , I) is a generalized polygon. Ampak v stavbe se ne bi poglabljal.

\begin{trditev}[Projektivne ravnine so natanko posplošeni trikotniki]
    Dokažimo ekvivalenco deficij takole:

    Vzemimo posplošen trikotnik $\Gamma = (\PP, \LL, I)$, ter dokažimo, 
    da je $\Gamma$ projektivna ravnina. Ker velja $n=3$, je ožina grafa 
    $G$ incidenčnega grafa $\Gamma$ enaka 6, premer pa 3.
    
    Če sta $P_1$ in $P_2$ poljubni točki, je razdalja med njima v dvodelnem grafu
    vedno soda. Ker je premer grafa enak 3, je razdalja med $P_1$ in $P_2$ 
    natanko 2, torej obstaja črta $l$, ki je v relaciji z obema točkama. Če
    bi obstajala še druga črta $l'$, ki je v relaciji z obema točkama, potem bi
    v grafu $G$ obstajal cikel $C_4$, kar je v nasprotju z ožino $g=6$. 
    S tem smo pokazali, da je aksiom 1. iz definicije projektivne ravnine zadovoljen.

    Z enakem argumentom lahko pokažemo, da je tudi aksiom 2. iz 
    definicije projektivne ravnine zadovoljen. Tu vidimo tudi, da je red 
    vsake projektivne ravnine enak $(s, s)$, saj se točke in premice obnašajo 
    v enakem številu, ter da je stopnja vsakega vozlišča enaka $s+1$.

    Ker je stopnja vsakega vozlišča v $G$ vsaj 3, je tudi
     aksiom 3. iz definicije projektivne ravnine izpolnjen, namreč (TODO: najlažje 
     je razumeti vizualno).

    Obratno, vzemimo projektivno ravnino $\Gamma = (\PP, \LL, I)$, ter dokažimo,
    da je $\Gamma$ posplošen trikotnik. Ker je $\Gamma$ projektivna ravnina, velja, 
    da je ožina njenega incidenčnega grafa ... TODO.



\end{trditev}






\section{TODO:}
Moorejev -> Mooreov.
Dokončaj izpeljavo pogojev. Prikaži izpeljavo za Hoffman-Singeltonov graf.
Pokaži še kaj o problemu degree/girth, in kako vsak Moorejov graf dosega 
tudi tisto spodnjo mejo, ni pa vsak graf ki dosega mejo za girth/degree Moorejov.
Pokaži da greedy approach za generiranje takih grafov spodleti (Že pri hoffman-singeltonu), 
pri petersonu pa se da.

Dokaži da je Hoffmanov graf (50,7,0,1) krepko regularen. Definicija 
krepke regularnosti. Opis razdaljno regularnih grafov in lastnosti.

Na kratko še o "blizu" moorejevih grafih. Nato pa skok na generalizirane večkotnike 
za nekaj časa.

Povezava med problemom ožine ter problemom premera. Teoretični uvod 
v geometrije. Nekaj o Fanojevi geometriji PG(2, 2) in konstrukciji 
hoffman-ovega grafa s tem. Opis povezav med problemom ožine ter 
posplošenih večkotnikov

\begin{definicija}
Naj bo $G$ usmerjen multigraf. Naj bo $\Gamma$ abelova grupa.
Označimo z $D(G)$ množico usmerjenih polrobov (darts) grafa $G$.
Naj bo $\alpha : D(G) \to \Gamma$ funkcija, ki vsakemu polrobu priredi element iz $\Gamma$,
in naj velja
\[
\forall e \in D(G)\colon\quad \alpha(e^{-1}) = \alpha(e)^{-1}.
\]
Preslikavi $\alpha$ pravimo \emph{napetostna preslikava} (voltage assignment).

Dvig grafa $G$ z napetostno preslikavo $\alpha$ je graf $G^\alpha$, definiran kot:
\begin{enumerate}
    \item $V(G^\alpha) = V(G) \times \Gamma$,
    \item $E(G^\alpha) = \{\, \bigl( (u,g), (v,\, g \cdot \alpha(e)) \bigr) \;\mid\; 
           e = uv \in D(G),\ g \in \Gamma \,\}$.
\end{enumerate}
\end{definicija}


\section{Working on:}

% Dynamic cage survey ima lepo definicijo, začni z definicijo kletke in mejo za vozlišča, 
% in potem tisti ki jo dosežejo so moorevoi, označi jo z M(k=degree, g), no ampak tak graf
% je regularen po definiciji, bi bilo treba predrugačiti trditev da itak so regularni tudi drugače
% . Vsi mooreovi grafi so kletke. minimalen (k, g)-graf je kletka.

% Regarding (v), it should be noted that the only known symmetric generalized n-gons
% have prime power order.

% To this day, Sachs’ result is the only constructive proof of the existence of (k, g)-graphs
% for any set of parameters k > 2, g > 3 that is completely non-algebraic.

% tista 5kotna zvezda se imenuje pentagram.

% Če želiš konstruirati Moorev graf z ožino g=6, moraš konstruirati 
% posplošen trikotnik (n=3). Posplošen trikotnik pa je v resnici le 
% drugo ime za projektivno ravnino.

% Levi graf je drugo ime za incidenčni graf.
% Heawoodov graf je preprosto incidenčni (Levijev) graf Fanojeve ravnine (ki je projektivna ravnina).
% Matrika sosednosti za incidenčni graf je preprosto [[0, M]; [M^T, 0]], kjer je M incidenčna matrika
% geometrije.

% However their classification is
% very difficult, since, for instance, every projective plane is a generalized
% triangle. (Feigh in higmann)
% To je "sveti gral" teh konstrukcij. Najbolj znana je Tits-ov posplošeni šesterokotnik.

% ki jo je razvil Hans Van Maldeghem.




%%%%%%%%%%%%%%%
%%%%%%%%%%%%%%%%
%%%%%%%%%%%%%%%%% 

\end{document}